\documentclass[12pt]{article}
\usepackage[T1]{fontenc}
\usepackage[margin=1in]{geometry}
\usepackage{amsmath,amssymb}
\usepackage{graphicx}
\usepackage{array}
\usepackage{booktabs}
\usepackage{threeparttable}
\usepackage{caption}
\usepackage{authblk}
\usepackage[authoryear,round]{natbib}
\usepackage{xurl}
\usepackage[hidelinks]{hyperref}
\usepackage{microtype}
\emergencystretch=1.5em
\renewcommand{\thesection}{\Roman{section}}
\renewcommand{\thesubsection}{\thesection.\Alph{subsection}}
% Possessive citation: one macro so a bibliography change cannot desynchronize
% hand-rolled \citeauthor{X}'s (\citeyear{X}) constructions.
\newcommand{\citepos}[1]{\citeauthor{#1}'s (\citeyear{#1})}

\title{The Securitization Trade-Off: Household Choice, Loan-Level Survival, and the Shortfall of Quantitative Tightening}
\author{Eugene Ong}
\affil{Carnegie Mellon University}
\date{}
\begin{document}
\maketitle
\begin{abstract}
\noindent The Federal Reserve's Quantitative Tightening (QT) program relied on passive runoff from its agency mortgage-backed securities (MBS) portfolio. The U.S. par-payoff mortgage structure restricted prepayments during the 2022--2025 rate shock, and MBS runoff fell \$764.7 billion short of the phased redemption caps. We quantify the shortfall with a 10,000-household agent-based model and a loan-level hazard framework built from Freddie Mac performance data. Synthetic-population household choice recovers 11.9\% of the benchmark. On the accounting basis consistent with the benchmark, loan-level survival recovers 97.9\% in its literature-calibrated microsimulation form (92.2\%--99.2\% across covariate calibrations; composed and estimated-form variants in the introduction). A no-lock-in null recovers 88.7\% on the same basis: the lock-in elasticity's identified contribution is a +9.2-percentage-point margin --- comparable in magnitude to the covariate-calibration sensitivity --- and a monotone calibration-band response. A Danish market-value-repurchase counterfactual, recalibrated to estimated Danish elasticities, collapses an initially large institutional gap by an order of magnitude, to a gap of roughly \$100 billion in magnitude whose sign is not robustly identified. TBA-market liquidity is purchased with household mobility; we find no robust institutional cash-flow leg to that trade-off.
\end{abstract}

\vspace{0.5\baselineskip}
{\footnotesize\noindent\textbf{Keywords:} mortgage prepayment; lock-in effect; quantitative tightening; mortgage-backed securities; agent-based models; hazard models.\\
\textbf{JEL codes:} E52, E58, G21, R31.\par}

\section{Introduction}\label{sec:intro}

For most of its history, the Federal Reserve operated under a ``Treasuries-only'' framework, using sovereign debt with fixed, predictable maturities to conduct monetary policy. During the 2008 Global Financial Crisis, and again during the 2020 COVID-19 pandemic, the central bank departed from this framework, launching aggressive Quantitative Easing (QE) programs that purchased trillions of dollars in Agency MBS to stimulate the consumer economy directly.

This intervention suppressed 30-year mortgage rates to historic lows of roughly 2.6\%, triggering a large refinance boom. As inflation surged in 2022, the Federal Reserve attempted to reverse this stimulus through Quantitative Tightening (QT), establishing a passive roll-off target intended to drain up to \$35 billion per month from its SOMA portfolio as homeowners relocated naturally --- a cap phased in from \$17.5 billion per month, stepping up to \$35 billion three months later.

The central argument of this paper is that the resulting QT shortfall is not incidental, but the predictable consequence of a specific institutional design choice we term the Securitization Trade-Off, and we use the term in a precise, two-sided sense throughout: the pool homogeneity that generates TBA-market liquidity is purchased at the price of household mobility. A third design objective --- extension-protected institutional cash flow --- is often framed as completing a trilemma; Section~\ref{sec:conclusion} presents our evidence that this third leg does not trade off sharply against mobility, so the operative exchange is the two-sided one. U.S. policymakers built the Agency MBS market around the To-Be-Announced (TBA) forward market, which requires deliverable pools to be homogeneous and therefore requires the ``par-payoff'' rule under which a homeowner must extinguish an existing mortgage at full face value to relocate. That same homogenization is what makes continuous household mobility and predictable central bank cash flow difficult to guarantee simultaneously. Section~\ref{sec:conclusion} returns to this framework in light of the evidence. Section~\ref{sec:abm}'s Danish counterfactual initially implied a \$925.5 billion institutional gap; recalibrating against separately-estimated Danish elasticities (Section~\ref{sec:abm-danish}) collapses this by an order of magnitude, to a gap of roughly \$100 billion in magnitude under the shared-accounting hybrid pipeline (13\% of the benchmark; point estimate $-\$99.9$ billion), and a sensitivity sweep shows the sign is not robustly identified --- a correction we treat as informative about where the trade-off actually binds rather than as a setback.

Unlike standard U.S. Treasuries, U.S. mortgages carry asymmetric, state-dependent payoff profiles. As the Federal Reserve's policy tightening pushed the 30-year fixed mortgage rate past 7.0\% --- a rise decomposable into the 10-year Treasury yield's climb and a widening of the primary-mortgage-to-10-year spread from roughly 170 to over 300 basis points at its 2022--2023 peak (FRED series MORTGAGE30US less DGS10); the decomposition of that spread into MBS supply, prepayment-risk repricing, and convexity-hedging components follows \citet{boyarchenko2019}, whose framework predates and does not document the 2022--2023 magnitudes --- the par-payoff rule imposed a substantial financial penalty on relocation, producing a ``lock-in effect'' that suppressed prepayment activity across the U.S. housing market. Because mortgages did not prepay, expected principal cash flows from the Federal Reserve's MBS portfolio fell far below schedule, extending the portfolio's effective maturity. This extension prevented the central bank from meeting its monthly QT target, stranding an estimated \$764.7 billion in expected SOMA cash flow over the tightening cycle, measured against the redemption cap (methodology and interpretation in Section~\ref{sec:method-benchmark}).

Standard macroeconomic models often underestimate the severity of this extension risk because they assume frictionless, rational household behavior. In practice, the lock-in effect was compounded by historic housing inventory shortages and depressed consumer sentiment.

This paper quantifies the mechanics of this shortfall using two structurally distinct estimators. The first is a 10,000-household Agent-Based Model parameterized with dynamic macroeconomic frictions, used to estimate the magnitude of stranded cash flows and to decompose that estimate against the empirical benchmark. The second, introduced after the ABM's residual proved persistent and unresponsive to additional behavioral realism, is a reduced-form hazard framework estimated on loan-level survival data that abandons the household utility function altogether. As Sections~\ref{sec:abm} and~\ref{sec:hazard} show, synthetic-population household choice explains only a small share of the shortfall (11.9\%; 10--17\% across seed noise), while the loan-level hazard framework recovers nearly all of it at the aggregate level --- 97.9\% in its literature-calibrated microsimulation form on the benchmark-consistent shared accounting basis (92.2\%--99.2\% across covariate calibrations), 100.0\% (approximately 94.3\%--101.3\% across those calibrations) once the sample's coupon mix is reweighted to the SOMA book under the ex-ante priors, and 110.6\% in its estimated cohort-hazard form, an aggregate-level figure only, since its forward simulation carries no useful propagated precision (Section~\ref{sec:patha}); the standalone-scorer figures reported alongside (107.0\% and 119.7\%) differ from these by a pure accounting layer of 9.1 percentage points --- with path-level diagnostics that support one of its two specifications considerably more strongly than the other; that hazard framework, not the ABM, is this paper's central empirical contribution. (Section~\ref{sec:robustness} confirms that the ABM's population size is not the source of this gap.) A counterfactual simulation based on the Danish market-value repurchase regime further shows that systemic duration risk is, at least in part, a consequence of U.S. institutional design rather than solely an unavoidable feature of high interest rates.

\section{Literature Review}\label{sec:lit}

The theoretical foundation of mortgage prepayment assumes ``ruthless exercise,'' wherein borrowers refinance instantaneously as market conditions dictate \citep{dunn1981}. Subsequent literature recognized behavioral frictions: \citet{schwartz1989} estimate an empirical prepayment function from aggregate GNMA experience with maximum-likelihood hazard techniques and embed it in a valuation framework --- the econometric ancestor of the loan-level survival approach we adopt in Section~\ref{sec:hazard} --- while \citet{richard1989} formalize the industry practice of modeling prepayment as an empirical function of the refinancing incentive, seasoning, seasonality, and a pool's own prepayment history. \citet{stanton1995} introduced heterogeneous transaction costs, smoothing the instantaneous refinancing cliff into a predictable S-curve, and \citet{quigley1987,quigley2002} established the lock-in effect as a persistent friction limiting mobility during rising-rate environments.

The severity of the 2022--2025 rate shock generated substantial empirical literature documenting the breakdown of standard prepayment assumptions. \citet{fonseca2024} documented the steep decline of U.S. mobility and refinancing activity, while \citet{batzer2024} estimated that U.S. households faced roughly \$2.4 trillion in foregone capital gains had they chosen to relocate. \citet{liebersohn2024} quantify the same lock-in mechanism in hazard form, estimating that each percentage point of prevailing rate above a borrower's locked-in coupon reduces quarterly mobility probability by 5.5\% to 7.7\% depending on specification --- the elasticity band our Path B microsimulation adopts as its calibration (Section~\ref{sec:pathb}).

A related strand of literature decomposes how much of the observed mobility decline a rational, rate-gap-driven framework can explain, leaving the remainder as an unexplained residual. \citet{aladangady2024} construct a counterfactual rate-gap distribution holding the 2021 rate environment fixed and find that rate-gap-driven lock-in explains only 44\% of the actual decline in mortgage borrower mobility between 2021 and 2022, after correcting for selection bias from forgone refinancing opportunities. This decomposition approach --- comparing a model's prediction against the observed empirical outcome to isolate an unexplained residual --- is methodologically analogous to the approach we take in Sections~\ref{sec:abm} and~\ref{sec:hazard}, though we apply it to dynamically modeled CPR and cash-flow shortfalls rather than to mobility rates directly, and deploy two structurally distinct estimators rather than one. \citet{gerardi2024} similarly compare a structural search-and-matching model of lock-in against empirical lifecycle migration patterns, finding that their model accounts for most, but not all, of the observed disruption to housing market liquidity.

A separate behavioral finance literature provides theoretical grounding for what such residuals likely represent. \citet{keys2016} document that more than 20\% of U.S. borrowers fail to refinance even when doing so would be unambiguously in their financial interest, attributing this to procrastination and variation in financial sophistication. \citet{andersen2020} reach a similar conclusion using Danish administrative data, finding that inattention and inertia, rather than transaction costs, account for a substantial share of household failure to act on refinancing opportunities. \citet{andersen2022} further show, using Danish housing data, that households exhibit reference dependence and loss aversion relative to their original purchase price, weighting losses roughly 2.5 times more heavily than equivalent gains. \citet{johnson2019} examine sluggish refinancing using administrative, survey, and field-experimental data, finding that among the behavioral factors they test, suspicion of lenders' motives --- rather than standard rational frictions such as time preference or risk aversion --- is the factor most consistently related to households' failure to accept profitable refinancing offers, lending further support to a behavioral interpretation of the residual. \citet{bissiri2014} formalize this distinction directly, defining ``behavioral risk'' as the share of prepayment decisions not driven by a fully rational financial criterion, an operationalization more recently applied by \citet{perotti2024} to model behavioral uncertainty as a stochastic overlay on an otherwise rational prepayment framework.

A further strand treats prepayment as a loan-level survival process shaped by pool composition and servicer behavior rather than by borrower optimization. Prepayment ``burnout'' --- the tendency for a pool's aggregate refinancing sensitivity to decline over time as its most rate-sensitive borrowers prepay earlier and self-select out of the remaining pool --- is a standard feature of industry MBS valuation models --- it is one of the four factors of \citepos{richard1989} canonical prepayment function, alongside the Public Securities Association seasoning convention --- but is difficult to reconcile with a fixed, representative-agent population. Our hazard framework, introduced in Section~\ref{sec:hazard}, is designed specifically to accommodate this loan-level survival structure.

Fixed-income literature separately highlights the systemic risks generated by the lock-in effect. \citet{jiang2023} documented that duration extension across bank asset portfolios, prominently including low-coupon MBS, generated roughly \$2 trillion in unrecognized mark-to-market losses across the U.S. banking sector, contributing to self-fulfilling solvency runs including the collapse of Silicon Valley Bank. \citet{smith2023} demonstrated that QT relies heavily on predictable MBS runoff schedules, which are structurally broken by these duration extensions.

The Danish mortgage securitization system offers a tested structural alternative to negative convexity. \citet{berg2018} outlined the Danish ``Balance Principle,'' which requires strict one-to-one match funding. \citet{berger2026} demonstrated that under these rules, household mobility remains insulated from rising rates because borrowers hold the right to repurchase their outstanding mortgages at market value, neutralizing extension risk for the bondholder.

A final strand situates the lock-in effect within the broader theory of monetary policy transmission. \citet{berger2021} demonstrate that the Federal Reserve's ability to stimulate the economy through rate cuts is path-dependent, because prepayment behavior depends on the outstanding distribution of in-the-money refinancing incentives. This paper extends that logic to the tightening direction: the same path dependence that weakens easing after a low-rate period also weakens tightening after a low-rate period, since the resulting stock of deeply in-the-money, rate-locked mortgages cannot prepay regardless of how aggressively the central bank subsequently raises rates. \citet{eichenbaum2022} formalize this state dependence for the easing direction --- the potency of monetary policy depends on the distribution of outstanding rate gaps --- and the tightening-side friction studied here is that mechanism's mirror image, operating on the rate-gap distribution the 2020--2021 refinancing wave left behind. \citet{beraja2019} provide complementary regional evidence that the strength of the refinancing channel varies systematically with the local distribution of in-the-money mortgages. \citet{vickery2013} document that the liquidity premium of the TBA market derives specifically from the homogeneity of deliverable pools, finding that TBA-eligible securities trade at significantly lower cost than otherwise similar specified pools; this homogeneity requirement is the institutional constraint that, in Section~\ref{sec:conclusion}, we argue makes the Danish market-value repurchase structure difficult to reconcile with the current U.S. secondary market absent a broader redesign of agency securitization rules.

\section{Empirical Methodology}\label{sec:method}

We quantify the Federal Reserve's QT shortfall using two structurally distinct estimators, introduced together here and each developed in full over the remainder of the paper.

\subsection{Two Structurally Distinct Estimators}\label{sec:method-estimators}

The first estimator is a 10,000-household Agent-Based Model (ABM), specified in Section~\ref{sec:method-abm}, which draws a synthetic population calibrated to aggregate FRED medians and tests whether a utility-maximizing household decision rule --- rational or behaviorally extended --- can account for the shortfall; its results are reported in Section~\ref{sec:abm}. The second is a hazard framework, specified in full in Section~\ref{sec:hazard}, which abandons the household utility function entirely and models loan-level survival directly on real Freddie Mac loan-level data: a stratum-month discrete-time hazard model estimated by Poisson pseudo-maximum likelihood on the full 2017--2021 origination universe aggregated to stratum-months (Path A), and a loan-level microsimulation with competing prepayment and default risks on a stratified 75,000-loan sample drawn from that universe (Path B).

Both estimators are evaluated against the same \$764.7 billion empirical benchmark (Section~\ref{sec:method-benchmark}) and the same June 2022--November 2025 QT window, so their results are directly comparable despite drawing on different data and different behavioral assumptions. The two estimators differ in modeling paradigm --- household choice versus loan-level survival --- and, on the hazard side, in the use of real rather than synthetic loan-level data; Section~\ref{sec:robustness} shows that the difference in population size between them (10,000 households versus Path B's 75,000 simulated loans, with Path A estimated on the full origination universe) is not what drives the gap in their explanatory power, though the data-source difference remains an unresolved confound that qualifies the paradigm interpretation throughout.

\subsection{Agent-Based Model: Specification}\label{sec:method-abm}

We adopt an agent-based approach rather than a representative-agent framework because the phenomenon under study is fundamentally heterogeneous: households differ in coupon, loan age, and exposure to search and sentiment frictions, and these differences interact non-linearly with the level of interest rates. \citet{geanakoplos2012} argue for, and demonstrate the feasibility of, agent-based models of mortgage prepayment as an alternative to representative-agent approaches in environments where behavioral and search frictions vary across the population and over time, a property directly relevant here since the frictions we model (housing search costs and consumer sentiment) are themselves time-varying. The broader methodological case for agent-based computational economics as an alternative to representative-agent modeling in settings with meaningful heterogeneity is set out in \citet{tesfatsion2006}.

The ABM simulates a population of N = 10,000 heterogeneous households. Initial agent endowments (income, home value) are drawn from log-normal distributions calibrated to Federal Reserve Economic Data (FRED) medians. In the production multi-vintage specification of Section~\ref{sec:method-frictions}, each household's mortgage coupon and origination date are drawn from the Federal Reserve's actual holdings composition --- the live CUSIP-level SOMA parse of Section~\ref{sec:abm-results}, across seven coupon cohorts. A single-cohort variant, used for exposition and mechanism analysis, instead endows every agent with a 30-year fixed-rate mortgage at a 3.0\% coupon aged 60 months at the simulation start; where the text below discusses a single representative coupon, it refers to this pedagogical case, not to the production population.

Population size is not what separates the ABM from the hazard framework: Section~\ref{sec:robustness} reports a scale-sensitivity replication confirming that the gap in explanatory power is not an artifact of sample size; whether it reflects modeling paradigm or data source is taken up in Section~\ref{sec:robustness-crossdesign}.

The decision to relocate is governed by comparing the utility of moving against the aggregate financial penalty of doing so. The monthly payment is calculated via the standard annuity formula
\begin{equation}\label{eq:annuity}
M = P\,\frac{r\,(1+r)^{n}}{(1+r)^{n}-1},
\end{equation}
where $P$ is the outstanding principal, $r$ is the monthly interest rate, and $n$ is the remaining term in months. For computational tractability across a three-dimensional grid of market rate, dynamic friction, and rate velocity (and, in the multi-vintage specification, across seven coupon cohorts), the population's move/stay decisions are evaluated household-by-household at each grid point and aggregated to a pool-level annualized Conditional Prepayment Rate (CPR) surface --- CPR is a pool-level convention; the household-level object is the monthly move probability --- which is precomputed once and interpolated month-by-month against the live population of simulated market conditions, rather than re-solved via per-month household loops. This reduces the cost of the full grid by roughly two orders of magnitude and makes fifty-seed Monte Carlo estimation feasible.

The simulation contrasts two institutional payoff rules. Under U.S. par-payoff, the cost to retire the mortgage is strictly equal to the outstanding principal --- the negatively convex case, in which the investor's upside is capped while extension risk is bounded only by the scheduled-amortization profile: a 14.7-year weighted-average life at zero CPR against the no-shock 3.4-year life (Table~\ref{tab:wal}), an order-of-magnitude extension ceiling rather than an unbounded one. Under Danish market-value repurchase, the payoff cost is the present value of the remaining scheduled payments discounted at the current 30-year primary mortgage rate (FRED MORTGAGE30US), capped at par. Two qualifications attach to this proxy. First, it is an option-free present value discounted at the primary rate: it ignores the embedded par call and any option-adjusted spread, and the primary rate embeds the primary--secondary spread (guarantee fee, servicing, originator margin), where the bond-market discounting a true buyback would use sits below it --- both choices overstate the borrower's buyback discount as rates rise. These biases inflate the mechanism-substitution Danish CPR of Section~\ref{sec:abm-danish} (the since-superseded 47.1\%) but are exactly immaterial to the Berger-recalibrated production finding: that leg's prepayment hazard is the imported flat moving elasticity, not this pricing proxy, and a replication-package diagnostic that re-runs the Danish leg with the discount rate shifted down by 50--100 basis points moves trapped liquidity by \$0.00 billion, while flipping 7,803--32,486 of the 1.68 million loan-month present-value classifications that no downstream quantity in that leg consumes. Second, what the Danish rule neutralizes is extension risk for the bondholder \citep{berger2026}: the borrower retains the par call, so the rate-decline side keeps its price cap, and we use ``convexity-neutral'' below strictly in this rate-rise sense. Under the Danish rule, as market rates rise above 3.0\%, the bond's present value falls, allowing the borrower to extinguish the debt at a discount and guaranteeing continuous cash flow generation. The Danish counterfactual balance is simulated forward on a dynamic, declining-balance basis, so that each month's elevated Danish roll-off is applied to the already-shrunk balance carried over from the prior month and the simulated payoff path compounds correctly over the multi-year QT window.

\subsection{Macroeconomic Friction and Structural Extensions}\label{sec:method-frictions}

To quantify the behavioral friction gap in standard models, we scale base transaction friction using two FRED series: (1) search friction, measured as the shortfall in Active Listing Count (ACTLISCOUUS) against a 2017--2019 baseline --- the last stable pre-pandemic period before the 2020/21 rate-driven frenzy and the 2022 inventory collapse; and (2) psychological friction, indexed from University of Michigan Consumer Sentiment (UMCSENT) against a baseline of 85.0, with readings below that line scored as incremental friction at +5bp per point short, capped at +150bp during acute sentiment stress. This produces a dynamic friction parameter --- a one-time moving transaction cost expressed as a share of the household's current home value --- that scales between 8.23\% and 10.26\% of home value over the sample.

Three literature-grounded behavioral extensions were added to the household decision function, each parameterized entirely from external sources with no parameter tuned to match the empirical benchmark: a debt-to-income (DTI) constraint that vetoes any move in which the new housing payment would exceed 43\% of gross monthly income --- a front-end heuristic motivated by the pre-2021 General QM debt-to-income limit (12 C.F.R. \S{} 1026.43(e)(2)(vi)); the regulatory threshold was a back-end total-debt ratio and was superseded by the price-based General QM definition in 2021, so we use it as a behavioral anchor rather than a live regulatory constraint, and the model carries no total-debt data against which to test the threshold (Section~\ref{sec:limitations}); asymmetric loss aversion that scales perceived payment increases by 2.25$\times$ while leaving decreases unscaled (loss aversion: \citealp{kahneman1979}; the 2.25 coefficient: \citealp{tversky1992} --- \citepos{andersen2022} housing-market estimate of roughly 2.5 is close, and because any coefficient in this range triggers the same turnover-floor recalibration described below, the artifact interpretation of Section~\ref{sec:abm-results} does not hinge on the choice); and a wait-and-see freeze that holds 20\% of otherwise-mobile households in place whenever the trailing six-month change in mortgage rates exceeds 150bp, modeled as a fixed per-agent trait so the CPR surface remains deterministic.

Because loss aversion steepens the penalty function facing every household, and the model's mobility-desire parameter is independently calibrated to reproduce a real-world involuntary-turnover floor of 4--5\% CPR at 8\% market rates --- the turnover speed observed on deeply out-of-the-money discount cohorts during 2023--2024, when refinancing incentives were absent and remaining prepayment reflected death, divorce, and forced relocation, we hold that empirical floor fixed and re-run the calibration under the steeper penalty function, which requires nearly tripling the mobility-desire scale parameter. This recalibration is a mechanical consequence of holding a real-world anchor fixed while the penalty function around it changes, and is the principal driver of the results in Section~\ref{sec:abm-results}.

We state the underlying rule explicitly, because it recurs: whenever a specification change alters the penalty or covariate structure facing every household, the mobility-desire scale --- and only that parameter --- is re-derived so that the simulation continues to reproduce the empirically observed 4--5\% involuntary-turnover floor; every other parameter is left untouched. The same rule generates the recalibrated variant of the cross-design test (Section~\ref{sec:robustness-crossdesign}), whose frozen variant is defined precisely by suspending it, and Section~\ref{sec:robustness-extensions} reports how much work this single choice does.

An income-scaled curtailment channel was added as an extension specified ex ante: a diagnostic correlation between the ABM's monthly CPR over-prediction and year-over-year growth in real disposable income (FRED series DSPIC96) of $r = -0.487$, concentrated in the worst over-prediction months (June--September 2022), led us to predict ex ante that curtailment --- a strictly additive channel --- would worsen rather than improve the aggregate match. Curtailment was implemented as an income-scaled addition to simulated roll-off, mapping real income growth linearly from zero curtailment at $-2$\% growth to a full-scale annualized rate of 1.2\% CPR at +4\% growth, consistent with GSE daily prepayment report estimates.

Finally, the single representative 3.0\% coupon assumed elsewhere was replaced, as a diagnostic, with a multi-vintage composition: SOMA holdings bucketed into seven coupon cohorts yield a weighted-average coupon of approximately 2.55\% on the 30-year book, with the 2.0\% and 2.5\% buckets accounting for roughly 69\% of 30-year face value. Because a lower coupon implies a larger gap to prevailing market rates, our ex ante hypothesis was that this composition would increase predicted trapped liquidity toward the empirical benchmark. Fifteen-year MBS (approximately 9\% of SOMA face value) are excluded from this pass, and cohort origination dates are approximated from parsed maturity strings under a 360-month term assumption. Section~\ref{sec:abm-results} describes the production replacement of this diagnostic with a live, CUSIP-level pull of SOMA holdings.

\subsection{Empirical Validation Benchmark}\label{sec:method-benchmark}

We construct the empirical shortfall benchmark from the New York Fed's SOMA current face value series, netting realized roll-off against the actual phased redemption cap schedule --- \$17.5 billion per month from June through August 2022, stepping up to \$35 billion per month from September 2022 onward --- strictly within the active QT window of June 2022 through November 2025. Net accumulation, rather than one-sided accumulation of shortfall months only, is essential for comparability with the Danish counterfactual of Section~\ref{sec:abm-danish}, where simulated roll-off routinely exceeds the cap.

This methodology yields a validated empirical shortfall of approximately \$764.7 billion (unrounded: \$764.748 billion; every percentage-of-benchmark figure in this paper divides by this value), which we adopt as the benchmark against which we evaluate the ABM's and the hazard framework's simulated trapped liquidity throughout Sections~\ref{sec:abm} and~\ref{sec:hazard}. The SOMA face underlying this benchmark spans the Federal Reserve's full agency MBS book --- including Ginnie Mae collateral, 20.4\% of face --- and 66.2\% of that face lies in the 2017--2021 vintages covered by the Freddie Mac loan-level sample of Section~\ref{sec:hazard}; the hazard framework's loan-level foundation therefore directly samples neither Ginnie Mae collateral nor roughly one-third of the book by vintage. The empirical CPR back-out nets cohort-weighted scheduled amortization and an income-scaled curtailment assumption averaging approximately 0.84\% annualized, consistent with GSE daily prepayment report estimates. Two implementations of this back-out exist in the pipeline --- the ABM side nets cohort-weighted, term-aware scheduled amortization, while the hazard side uses a single-pool amortization approximation --- and any correlation computed against them differs only at the third decimal; Table~\ref{tab:estimators}'s note states which rows use which. A full robustness sweep across data source and cap-schedule assumptions is reported in Section~\ref{sec:robustness}.

A conceptual caveat attaches to the interpretation of this figure, and we address it directly rather than leave it implicit. The benchmark measures shortfall relative to the redemption cap, and the cap is a ceiling on runoff rather than a forecast of it: market participants understood relatively early in the tightening cycle that MBS runoff was unlikely to reach the \$35 billion cap in most months, so the shortfall relative to the cap is an upper bound on the shortfall relative to any reasonable ex ante expectation. We nonetheless adopt the cap as the counterfactual for two reasons. First, the cap is the only policy-stated quantity available: the Federal Reserve published no expected MBS runoff path, so any alternative counterfactual would require a model of expected prepayment --- precisely the object under study --- and would render the benchmark circular with respect to the estimators evaluated against it. Second, the cap-relative construction is the policy-relevant quantity: it measures the gap between the maximum balance-sheet reduction the FOMC authorized and the reduction the mortgage structure delivered. Throughout the paper, ``trapped liquidity'' and ``shortfall'' accordingly denote the shortfall relative to the cap, and the percentages reported for each estimator are shares of that cap-relative quantity; because all three estimators are measured against the same construction, their comparative ranking is unaffected by this interpretive choice.

\section{Agent-Based Model Results}\label{sec:abm}

Synthetic-population household choice does not explain the Federal Reserve's trapped liquidity. This section establishes that result and the falsification tests behind it. Every specification tested --- rational baseline, behavioral extensions, and a burnout variant --- moves further from, or fails to close, the \$764.7 billion benchmark, so we report these results briefly and move directly to the structurally independent hazard framework of Section~\ref{sec:hazard}, this paper's central empirical contribution and the estimator that recovers, at the aggregate level, nearly all of the benchmark that the synthetic-population ABM cannot. Figure~\ref{fig:recovery} summarizes the estimator comparison developed in this section and the next.

\begin{figure}[t]
\centering
\includegraphics[width=\textwidth]{fig1_recovery_by_estimator}
\caption{Benchmark recovery by estimator, on the shared accounting basis of Section~\ref{sec:robustness-hybrid} (standalone-scorer values in parentheses; ABM-family rows run through that layer natively). Raw lag-0 correlations are annotated for reference only: no estimator's co-movement is positive once the window's shared trend is removed (Section~\ref{sec:pathb}), and level recovery alone does not rank models. Constructed from the frozen run artifacts in the replication package.}
\label{fig:recovery}
\end{figure}

\subsection{Baseline, Production Estimate, and Behavioral Extensions}\label{sec:abm-results}

A purely rational baseline projects a structural cash flow deficit of \$544.0 billion (71.1\% of benchmark). The production ABM builds on this baseline in four ways: scheduled amortization, a live CUSIP-level parse of SOMA holdings, a corrected empirical CPR back-out, and a settlement-lag kernel. A robustness test found that the kernel does not improve fit: with the kernel applied, the best cross-correlation is $r = +0.192$ at lag $-3$, indistinguishable from the $+0.195$ the kernel-free production run reaches at the same lag, so the production specification omits it. Together, these additions yield a production estimate of \$91.0 billion in trapped liquidity, or 11.9\% of the benchmark (the frozen seed-42 replication draw; the Monte Carlo distribution is reported below), following the incorporation of the previously excluded 15-year MBS book into cohort weighting and scheduled amortization (structural terms only; voluntary CPR for 15-year cohorts continues to use the same-coupon 30-year surface, since the payment-delta gate is unreliable for short-amortization loans --- see Section~\ref{sec:robustness-pipeline}). The CPR path fit is equally weak. Simulated mean CPR runs at 11.68\% against an empirical 5.14\%, and the raw path correlation is negative at lag 0 --- $r = -0.318$, $R^2 = -6.984$\footnote{Throughout, $R^2$ is computed as $1-\mathrm{SSE}/\mathrm{SST}$ against a mean-only null; negative values indicate fit worse than a constant at the empirical mean.} --- with no value inside a $\pm$3-month window exceeding the +0.19 at the three-month lead. (A fifty-seed Monte Carlo re-draw under the same specification has mean \$103.7 billion --- 13.6\% of benchmark --- with a seed-level standard deviation of \$24.5 billion, a $\pm$1~SD share band of 10.4--16.8\%; the 95\% CI of the mean, \$96.9--\$110.5 billion, reflects only simulation noise in the mean itself, not seed dispersion. The frozen draw sits near the 34th percentile of the seed distribution, a percentile itself estimated only to within roughly $\pm$7 percentage points from 50 draws, so the ABM's share is better read as on the order of 10--17\% than as a three-digit point estimate.)

Layering in three literature-grounded behavioral extensions moves the model to 54.9\%, then 45.1\%, then 33.7\% of the benchmark --- and each move is a further falsification, not progress toward it. The first extension (a debt-to-income constraint, asymmetric loss aversion, and a rate-shock freeze) lifts the estimate to 54.9\%, but only because holding a real-world turnover floor fixed under a steeper loss-aversion penalty forces a mechanical recalibration that inflates the model's rate sensitivity; it is an artifact of that recalibration, not evidence for loss aversion as a mechanism. The second, income-scaled curtailment, pulls the estimate down to 45.1\%, exactly as an ex ante income-growth diagnostic predicted. The third, multi-vintage coupon composition, pulls it down further to 33.7\% --- the opposite of our hypothesis that a wider coupon-rate gap should raise, not lower, predicted trapped liquidity. None of the three touches the CPR path's central failure: the contemporaneous correlation with the empirical series stays negative, and the best positive alignment anywhere in the $\pm$3-month window --- the $+0.19$ at the three-month lead noted above --- is weak and statistically indistinguishable from zero (a lag-3 coefficient has $42 - 3 = 39$ overlapping observations, so its i.i.d.\ reference band is $\pm 0.31$), since each extension operates on the level of simulated mobility rather than its timing. Full mechanism-level detail for each extension is reported in Section~\ref{sec:robustness-extensions}. Within the specification tested, added household-level realism does not recover the missing liquidity.

\subsection{Burnout Falsification}\label{sec:abm-burnout}

We tested whether path-dependent prepayment ``burnout,'' implemented as closed-population survivor selection, could explain the ABM's weak CPR fit. It failed outright: simulated CPR collapsed to approximately zero and trapped liquidity rose to 147\% of the benchmark, because a closed population of 10,000 agents depletes its mobile mass during the low-rate 2020--2021 window, before QT even begins --- an absorbing state incompatible with a real MBS pool, which continuously receives new originations and retains a baseline turnover rate. We do not use this specification in the production pipeline; we report it as a negative finding that motivates the cohort-level, rather than per-agent, treatment of burnout in Section~\ref{sec:hazard}'s hazard framework.

\subsection{Danish Counterfactual}\label{sec:abm-danish}

Sections~\ref{sec:abm-results} and~\ref{sec:abm-burnout} establish that household choice, under the U.S. par-payoff rule, does not explain the majority of the trapped liquidity --- so the natural next question is how much of that trap is a consequence of the par-payoff rule itself, rather than of household behavior under any institutional regime. The Danish market-value repurchase system, which severs the link between market rates and payoff cost, gives us a direct way to ask that question quantitatively rather than by assertion. Our first specification answered it by taking each framework's existing, U.S.-calibrated response function and feeding it a Danish-style payoff rule --- a mechanism substitution, not an estimated elasticity. That specification produced a Danish CPR of 47.1\% under the ABM and an institutional gap of approximately \$925.5 billion\footnote{\label{fn:manifest}All ABM headline figures cite frozen run \texttt{run-2026-07-04-15yr-foldin} (\path{abm/data/runs/run-2026-07-04-15yr-foldin/manifest.json}) in the replication package; the run is reproduced via \texttt{freeze\_run.py -{}-tag run-2026-07-04-15yr-foldin}. CPR means aggregate over the 42 active QT months (June 2022--November 2025). The institutional wedge is the mean over those months of the difference between the cohort-weighted Danish and U.S. simulated CPR series. The \$925.5 billion figure is retained in the text at first mention only to show the specification we started with before correcting it, not as a competing estimate; it is superseded by the Berger-recalibrated specification, \texttt{run-2026-07-05-berger}, reported below. Superseded figures from earlier drafts and specifications are catalogued, with provenance, in Appendix~\ref{app:ledger}.}. We report this only as the specification we started with; it is superseded by the recalibration below, which we adopt as the paper's actual finding on this question.

The mechanism-substitution approach conflates two behaviorally distinct channels, and \citet{berger2026} --- our source for the Danish institutional mechanism itself --- estimate them separately. The first is a moving channel: does removing the par-payoff penalty change how often households relocate? The second is a refinance-in-place channel: does the option to buy back debt at a market discount induce discount-buying behavior independent of moving? Danish administrative data show these behave very differently under a U.S. tax transplant. The moving channel is well-identified and close to flat with respect to the coupon gap --- Danish households move at an unconditional rate of 3.2\% per year regardless of how far their coupon sits below market, versus a U.S. moving-hazard slope of 0.57--1.20 percentage points per year per 100 basis points of coupon gap \citep{fonseca2024}. The refinance-in-place channel, by contrast, is substantially a Danish tax artifact: discount buybacks are attractive there because mortgage interest is deductible at a flat 33\% rate and realized capital gains on the buyback are tax-exempt. Under U.S. tax treatment (a 22\% mortgage-interest deduction and a 15\% capital-gains tax on the discount, per Berger et al.'s own counterfactual), this channel loses its shield, and their general-equilibrium counterfactual finds it shifts the equilibrium mortgage rate by only about one basis point --- economically negligible.

Recalibrating both frameworks against these two estimated channels, rather than against a single extrapolated mechanism, collapses the finding rather than reversing it with confidence. Importing Denmark's flat 3.2\% moving elasticity and treating the U.S.-transplant refinance-in-place channel as negligible (per Berger et al.'s own general-equilibrium result) yields a Danish CPR of approximately 3.4\% --- barely different from, and in some specifications below, the U.S. side. Under the paper's shared-accounting hybrid pipeline --- specified in Section~\ref{sec:robustness-hybrid}: the ABM's macro-accounting layer with Path B's microsimulated CPR as the micro-foundation --- which puts both regimes on the same empirically-grounded footing (U.S. 4.76\%, Danish 3.39\% mean CPR), the institutional gap becomes approximately $-\$99.9$ billion at this best-estimate input. Table~\ref{tab:danish} collects the three Danish specifications. But we do not report that figure as a reliable negative finding: Berger et al.'s general-equilibrium result puts the refinance-in-place channel near zero without pinning it to exactly zero, and a direct sensitivity sweep of that one input --- from 0\% up to the partial-equilibrium reduced-form estimate we otherwise rejected as an over-predicting ceiling ($\approx$18\% CPR) --- shows the hybrid pipeline's gap crosses from negative to positive at a contribution of just 1.4\% CPR, well within what a channel described only qualitatively as ``small'' could plausibly deliver. The ABM's own recalibrated run shows a larger, more sign-robust reversal ($-\$728.4$ billion, crossing zero only at a 12.6\% refinance-in-place contribution), but we do not treat this as reassuring either: the ABM's U.S.-side leg is itself over-predicted relative to the empirical benchmark (11.76\% simulated mean CPR under the Berger run's own calibration-stage freeze, versus 5.14\% empirical; cf.\ the production tag's 11.68\%, Section~\ref{sec:abm-results}), so comparing an inflated U.S. leg against an empirically-anchored Danish leg overstates whatever reversal it shows.


\begin{table}[t]
\centering
\small
\begin{threeparttable}
\caption{The Danish market-value-repurchase counterfactual under three specifications.}
\label{tab:danish}
\begin{tabular}{@{}lrrrr@{}}
\toprule
Specification & U.S. leg & Danish leg & Gap (US$-$DK) & Mean CPR US/DK \\
\midrule
Mechanism-extrapolated ABM\tnote{a} & $+\$91.0$B & $-\$834.5$B & $+\$925.5$B & 11.68\% / 47.10\% \\
Berger-recalibrated ABM\tnote{b} & $+\$84.5$B & $+\$812.9$B & $-\$728.4$B & 11.76\% / 3.40\% \\
Berger-recalibrated hybrid\tnote{c} & $+\$749.0$B & $+\$848.9$B & $-\$99.9$B & 4.76\% / 3.39\% \\
\bottomrule
\end{tabular}
\begin{tablenotes}[flushleft]\footnotesize
\item Sign convention: each leg is signed trapped liquidity --- the net shortfall of simulated roll-off relative to the phased redemption cap over the QT window --- so positive values mean runoff fell short of the cap (liquidity trapped) and negative values mean simulated runoff exceeded the phased cap path. Gap $=$ U.S. leg $-$ Danish leg; a negative gap means the Danish regime traps more than the U.S. regime.
\item[a] Superseded: applies the U.S.-calibrated mobility function to a Danish-style rate gap it was never estimated to describe (see text).
\item[b] Run \texttt{run-2026-07-05-berger}; its U.S. leg is re-simulated under that run's own calibration-stage freeze (native 15-year gate), not the \$91.0B production headline, so both legs and the gap share one frozen calibration. The Danish leg prepays more slowly under the tax-adjusted elasticities, trapping more than the U.S. leg.
\item[c] The paper's preferred estimate (shared-accounting hybrid pipeline, Section~\ref{sec:robustness-hybrid}); both legs are printed from the shared-layer scoring artifact (U.S. trapped \$749.0B = 97.9\% of benchmark; Danish trapped \$848.9B). Its sign is not robust --- the sweep of the refinance-in-place input crosses zero at a 1.4\% CPR contribution.
\end{tablenotes}
\end{threeparttable}
\end{table}

The original \$925.5 billion figure was, in retrospect, largely an artifact of applying a U.S.-calibrated mobility function to a Danish-style rate gap it was never estimated to describe --- extrapolation, not measurement. Once the two behavioral channels are estimated from the sources that actually study them, the institutional cash-flow case for the Danish alternative shrinks to something indistinguishable from zero, and we do not have a robust basis for signing it in either direction. This does not undo Sections~\ref{sec:abm} and~\ref{sec:hazard}'s central finding --- that loan-level and institutional plumbing, not household choice, explains the U.S. empirical shortfall --- since that comparison never depended on the Danish counterfactual. What it does undo is the implicit suggestion that a Danish-style payoff rule is the alterable fix, in either direction. Section~\ref{sec:conclusion} revises that framing accordingly.

\subsection{Interpretation}\label{sec:abm-interp}

Every ABM specification tested here either widens the gap to the benchmark or leaves it unchanged, and its CPR path stays uncorrelated or negatively correlated with the empirical path throughout. Population size is not the explanation --- Section~\ref{sec:robustness}'s scale-sensitivity replication rules that out --- so what remains is either that household-level decision-making does not govern the majority of the shortfall, or that this specification is missing a mechanism. Section~\ref{sec:hazard}'s hazard framework, built on an entirely different data source and modeling paradigm, tests the former directly. Whichever explanation holds, the ABM results already narrow the trade-off named in this paper's title: they weigh against household mobility choices, rational or behavioral, as the primary price the Federal Reserve paid for secondary-market liquidity, and point instead toward the loan-level and institutional channels Section~\ref{sec:hazard} takes up directly.

\section{The Hazard Framework: A Structurally Independent Estimator}\label{sec:hazard}

This is the paper's central empirical contribution. The ABM of Section~\ref{sec:abm} never approaches the hazard framework's recovery at the aggregate dollar level --- 97.9\% to 110.6\% of the benchmark on the shared accounting basis of Section~\ref{sec:robustness-hybrid}, 107.0\% to 119.7\% under the standalone scorers: its production specification recovers 11.9\%, and even its most inflated behavioral extension --- the 54.9\% figure we attribute in Section~\ref{sec:abm-results} to a mechanical recalibration artifact rather than to loss aversion as a mechanism --- falls well short. (The uncorrected rational baseline reaches 71.1\%, but that figure predates the scheduled-amortization, CUSIP-parse, and CPR-back-out corrections described in Section~\ref{sec:abm-results} and is not part of the production estimate family; we flag this explicitly so the baseline figure is not mistaken for a closer specification than the corrected ones that supersede it.) The hazard framework's recovery range, by contrast, uses an entirely different data source and modeling paradigm. We state at the outset the standard by which this success is, and is not, established: both hazard paths are evaluated against the same aggregate benchmark as the ABM, but only Path B's recovery is accurate on the benchmark-consistent accounting basis --- and no estimator's monthly co-movement with the empirical path, Path B's included, survives removal of the window's shared trend (Section~\ref{sec:pathb}). Path A's monthly path fit is no better than the ABM's --- a limitation Sections~\ref{sec:patha} and~\ref{sec:hazard-interp} quantify and weigh explicitly, rather than reserving for the limitations section, because it bears directly on how much evidential weight each recovery figure can carry.

\subsection{Motivation}\label{sec:hazard-motivation}

Section~\ref{sec:abm}'s pattern --- every structural correction and behavioral extension to the ABM either widens the gap to the empirical benchmark or leaves it unchanged, with a CPR path that remains uncorrelated or negatively correlated with the empirical series --- is difficult to reconcile with the hypothesis that the residual reflects a missing rational or behavioral household mechanism that a sufficiently rich utility function would eventually capture. It is more consistent with the hypothesis that the binding constraints on the empirical prepayment path lie at least partly outside the household decision function --- in loan-level heterogeneity, servicer pipeline dynamics, and cohort-level path dependence that a fixed, representative-household population struggles to represent.

This is the hypothesis the hazard framework introduced in Section~\ref{sec:method-estimators} is built to test directly. It abandons the household utility function entirely: the hazard framework contains no mobility-desire parameter, no loss-aversion coefficient, and no utility-maximizing decision rule; it instead models loan-level survival directly, using proportional hazards (proportional above an involuntary-turnover floor; Section~\ref{sec:pathb}) estimated on real loan-level data, with the closed-population survivor-selection failure of Section~\ref{sec:abm-burnout} directly motivating a cohort-level rather than per-loan treatment of burnout. The framework retains the same \$764.7 billion empirical benchmark and the same QT window definition as the ABM for comparability, and comprises two independent paths: a stratum-month discrete-time hazard model estimated on aggregated loan-level data, and a loan-level microsimulation calibrated to literature elasticities with explicit competing risks.

\subsection{Path A: Stratum-Month Discrete-Time Hazard Model}\label{sec:patha}

The first hazard path aggregates Freddie Mac loan-level performance data to the stratum-month level: strata $s$ are the 296 observed four-way cells of vintage $\times$ coupon $\times$ FICO bucket $\times$ LTV bucket, and the unit of observation is the pair $(s,t)$ of stratum and calendar month. The model is a discrete-time hazard with a log-exposure offset,
\begin{equation}\label{eq:pathA}
\log \mathbb{E}\!\left[y_{s,t} \mid x_{s,t}\right] = \log n_{s,t} + \alpha_{s} + f\!\left(a_{s,t}\right) + \beta_g\, g_{s,t} + \beta_b\, b_{s,t} + \beta_f\, \phi_t ,
\end{equation}
where $y_{s,t}$ is stratum $s$'s prepaid unpaid principal balance (UPB) in month $t$, in dollars, $n_{s,t}$ its UPB exposure, $x_{s,t}$ the right-hand-side covariates, $\alpha_{s}$ the stratum fixed effect (295 enter, one absorbed as reference), $f(\cdot)$ a loan-age seasoning spline in mean loan age $a_{s,t}$ (truncated quadratic power basis: a linear term plus $(a - k)_+^2$ at knots $k \in \{12, 24, 36, 60, 84, 120\}$ months, seven basis columns), $g_{s,t}$ the rate gap in 100bp units, $b_{s,t}$ the stratum burnout state, and $\phi_t$ the dynamic friction index. The outcome is dollars, not a count: with the log-exposure offset the modeled object is the prepayment rate $y_{s,t}/n_{s,t}$, a nonnegative fraction, so the model is estimated by Poisson pseudo-maximum likelihood, whose score equations remain consistent for the conditional mean without any distributional assumption on $y_{s,t}$; no dispersion statistic is meaningful for this pseudo-likelihood. All uncertainty statements below come from the two bootstrap schemes of Table~\ref{tab:bootstrap}: a stratum-level cluster bootstrap, robust to arbitrary within-stratum dependence and overdispersion but structurally unable to measure uncertainty in the friction coefficient --- $\phi_t$ varies only over time, so resampling strata reproduces the identical friction series in every replication --- and a temporal moving-block bootstrap that resamples the 36 training months in blocks of six with each month's full cross-section intact, which carries that burden (and, partially, burnout's). Table~\ref{tab:panel} reports the estimation panel's accounting.

\begin{table}[t]
\centering
\small
\begin{threeparttable}
\caption{Path A estimation panel.}
\label{tab:panel}
\begin{tabular}{@{}lrr@{}}
\toprule
 & Training & Holdout \\
\midrule
Performance months & Jan.\ 2021--Dec.\ 2023 & Jan.\ 2024--Sep.\ 2025 \\
Stratum-months & 10,176 & 6,077 \\
Prepaid UPB (PML outcome) & \$650.0B & \$131.3B \\
UPB exposure (month-sums) & \$72.5T & \$37.4T \\
\bottomrule
\end{tabular}
\begin{tablenotes}[flushleft]\footnotesize
\item The panel aggregates the full Freddie Mac 2017--2021 origination universe --- 8.84 million active loans at the window's peak month, not the 75,000-loan sample that Section~\ref{sec:pathb} simulates --- to 296 strata, and carries no sampling weights: exposure and the outcome are raw dollars, so the pseudo-maximum-likelihood fit is UPB-weighted by construction (a dollar of prepayment counts the same in any stratum-month). The temporal split at January 2024 is the ridge-selection holdout of the text; the performance window is January 2021 through September 2025.
\end{tablenotes}
\end{threeparttable}
\end{table}

The specification includes 295 stratum fixed effects --- 296 observed four-way combinations of vintage, coupon, FICO bucket, and LTV bucket, one absorbed as the reference --- which absorb each cohort's baseline prepayment speed so that the burnout term captures within-pool adverse selection over time rather than cross-sectional differences in pool quality. The burnout regressor is a within-stratum demeaned measure of the surviving loan stock, orthogonalized against the loan-age spline so that it is not mechanically collinear with seasoning. An earlier specification using four vintage-year dummies in place of the 295 stratum fixed effects returned a burnout coefficient of the wrong sign (+0.76); replacing the coarse vintage fixed effects with the full 295-stratum specification and demeaning burnout within stratum removes that wrong-signed artifact. The corrected estimate ($-0.13$) carries the theoretically expected sign --- a more heavily burned-out pool should prepay more slowly, holding age and rate gap fixed --- but the sign statement is a point-estimate statement only: the temporal bootstrap scheme of Table~\ref{tab:bootstrap} centers this coefficient on the opposite sign (mean $+1.23$, median $+0.52$), and under both schemes it is imprecise and not statistically distinguishable from zero, so we treat it as a control rather than a finding. Because 295 fixed effects make the model ill-conditioned under plain iteratively reweighted least squares --- 20 strata record no training prepayment event at all, so their fixed-effect maxima do not exist --- the fit carries a nominal ridge penalty $\lambda_{\mathrm{ridge}}$ on the full coefficient vector, with $\lambda_{\mathrm{ridge}}$ chosen over the grid $\{10^{-5}, 10^{-4}\}$ on a temporally blocked holdout comprising all stratum-months from January 2024 onward, so that no within-stratum information crosses the train/holdout split. We do not describe the penalty as what stabilizes the fit, because it is not: at these magnitudes the inestimable fixed effects are pinned by the warm-started optimizer rather than by the penalty --- a role verified against the unpenalized estimable-strata fit below --- and the penalty's operative role is as the fallback for resampled refits in the bootstrap machinery. Three facts about this device belong in print, all verified in the replication package. First, at these magnitudes the penalty is a numerical no-op on the full panel: the two grid points return bit-identical coefficient vectors from the IRLS warm start, so the grid selection is inconsequential and the production coefficients are the IRLS solution to machine precision --- ``mildly shrunk'' would overstate the penalty's effect, which is nil. Second, ridge on the full coefficient vector with one stratum absorbed as reference is not invariant to the choice of reference: in cold-start refits on resampled panels, swapping the reference stratum can move the standardized rate-gap coefficient by more than its own magnitude (an optimizer-path artifact of penalized refits in the presence of the 20 inestimable fixed effects; the bootstrap machinery below therefore warm-starts every replication and discards non-converged draws). Third, the production point itself is not an artifact of these choices: dropping the 20 zero-event strata (423 of 10,176 cells) and fitting unpenalized Poisson PML converges, is exactly invariant to the reference choice, and reproduces the production macro coefficients to within $0.002$ in standardized units --- rate gap $+0.6727$ against the production $+0.6734$ and friction $-0.0373$ against $-0.0371$ (third-decimal agreement), burnout $-0.1317$ against $-0.1301$, the larger burnout gap reflecting that dropping the 423 cells changes the burnout regressor's within-stratum demeaning and age-orthogonalization sample. Two honest caveats still attach: the same holdout later serves as the out-of-sample diagnostic, so the RMSE reported below is at least weakly favored by selection, and months never touched by any selection step do not exist in this design --- one reason Path A is held to a robustness role. The bootstrap reported below holds $\lambda_{\mathrm{ridge}}$ fixed rather than re-selecting it per replication.

Estimated on this panel, the specification yields a rate-gap coefficient of $+0.67$ (standardized; the training standard deviation of the rate gap is 221bp, so the equivalent per-100bp coefficient is $+0.30$), a burnout coefficient of $-0.13$, and a friction coefficient of $-0.037$, each with the theoretically expected sign. Point estimates alone overstate what is identified here, so Table~\ref{tab:bootstrap} reports cluster-bootstrap uncertainty for all three, including the bias diagnostics the bootstrap distribution itself demands: the rate-gap point estimate sits near the 3rd percentile of its own bootstrap distribution (bootstrap mean $+1.49$, median $+0.76$, estimated bias $+0.82$), a signature of resampling-induced upward skew under which percentile and basic intervals disagree --- the percentile interval excludes zero while the bias-reflecting basic interval does not. We therefore rest no claim on interval-based significance for any Path A coefficient. The defensible statements are about sign and imprecision: the rate-gap coefficient is positive in 197 of 198 converged replications (99.5\%), with a magnitude identified only up to a wide, construction-dependent interval; burnout and friction are not distinguishable from zero on any construction. All three enter the production simulation as fitted inputs, none as findings, and no result in this paper depends on their point values. Path A's aggregate recovery now carries the interval that pass produces: resimulating the QT window under each of the 198 converged replications (the three macro coefficients swapped in per draw; intercept, age spline, and fixed effects held at the production fit) yields a median of \$922.6 billion and an interquartile range of \$900.1--\$1{,}130.6 billion, but a 95\% percentile interval of $[-\$4{,}459.8, +\$1{,}190.1]$ billion: fifteen replications with extreme burnout or friction draws --- the two coefficients the bootstrap classifies as statistically zero --- destabilize the 42-month forward compounding entirely, and 83.8\% of draws land above the benchmark. Propagated coefficient uncertainty therefore leaves the forward simulation with no useful precision, which is an additional reason the \$915.0 billion figure is held to the aggregate-level, robustness role assigned below and is not carried in the abstract. Simulating forward --- draining each cohort's balance deterministically each month by balance times hazard, with fractional prepayment settling directly to simulated SOMA roll-off and no settlement-lag convolution --- this model recovers \$915.0 billion, or 119.7\%, of the \$764.7 billion empirical benchmark.

\begin{table}[t]
\centering
\footnotesize
\setlength{\tabcolsep}{4pt}
\begin{threeparttable}
\caption{Path A coefficient uncertainty under two bootstrap schemes.}
\label{tab:bootstrap}
\begin{tabular}{@{}lrrrrccr@{}}
\toprule
Coefficient & Point & Mean & Median & SE & Percentile CI & Basic CI & Draws $\le 0$ \\
\midrule
\multicolumn{8}{@{}l}{\textit{Stratum-cluster scheme (198 of 200 replications converged)}} \\
Rate gap (standardized) & $+0.673$ & $+1.491$ & $+0.763$ & $1.152$ & $[+0.60,\,+4.11]$ & $[-2.76,\,+0.74]$ & 0.5\% \\
Burnout (orthogonalized) & $-0.130$ & $-0.056$ & $-0.130$ & $1.363$ & $[-1.76,\,+1.43]$ & $[-1.69,\,+1.50]$ & 75.3\% \\
Friction & $-0.037$ & $+0.352$ & $+0.015$ & $0.506$ & $[-0.25,\,+1.38]$ & $[-1.46,\,+0.18]$ & 48.0\% \\
\midrule
\multicolumn{8}{@{}l}{\textit{Temporal moving-block scheme, block length 6 months (196 of 200 converged)}} \\
Rate gap (standardized) & $+0.673$ & $+3.313$ & $+2.556$ & $3.845$ & $[+0.44,\,+12.18]$ & $[-10.84,\,+0.91]$ & 0.5\% \\
Burnout (orthogonalized) & $-0.130$ & $+1.229$ & $+0.523$ & $2.586$ & $[-3.16,\,+6.16]$ & $[-6.42,\,+2.90]$ & 40.8\% \\
Friction & $-0.037$ & $+1.050$ & $+0.567$ & $2.802$ & $[-3.71,\,+7.69]$ & $[-7.76,\,+3.64]$ & 37.2\% \\
\bottomrule
\end{tabular}
\begin{tablenotes}[flushleft]\footnotesize
\item Stratum-cluster scheme: the observed vintage--coupon--FICO--LTV strata are resampled with replacement, each resampled stratum receives a fresh fixed-effect label so duplicated strata contribute independent fixed effects, and per-replication standardization (including the burnout orthogonalization, which is refit inside every replication) is rescaled to production standardized units. This scheme is structurally silent about the friction coefficient: the friction index varies only over time, so every replication sees the identical series, and the stratum-scheme friction row understates that coefficient's uncertainty by construction. Temporal scheme: the 36 training months are resampled in moving blocks of six with each month's cross-section kept intact, same refit and rescaling; it is the informative scheme for friction and, partially, burnout --- its friction SE (2.80) is roughly $5.5\times$ the stratum scheme's (0.51). The ridge penalty is held at the production $\lambda_{\mathrm{ridge}} = 10^{-4}$ rather than re-selected per replication (re-selection would bootstrap a different, coefficient-plus-model-selection estimator). Replications returning non-finite or absurd ($|\beta| > 20$) coefficients are counted as failures, not draws; the point refit reproduces the production coefficients to four decimals on the same 10,176 training cells. Coefficients are in production standardized units (training SD of the rate gap: 221bp). The stratum-scheme burnout median coincides with the point estimate at reported precision only (unrounded: median $-0.1297$, point $-0.1301$). Mean, median, and the implied bias (mean $-$ point) are reported because both schemes' rate-gap distributions are heavily right-skewed; under that skew the percentile and basic 95\% intervals disagree about zero, so interval-based significance is construction-dependent and the text relies on sign stability only --- which holds under both schemes: the rate-gap draw is positive in 99.5\% of converged replications under each. A BCa correction would require a 296-refit jackknife for the acceleration constant and is left to the replication package.
\end{tablenotes}
\end{threeparttable}
\end{table}

Three diagnostics temper this aggregate result, and we flag them here --- with the same prominence we gave the ABM's equivalent failures --- because they bear directly on how much weight the \$915.0 billion figure can carry. First, the holdout fit is weak in both weightings, and the units matter: the unweighted RMSE of approximately 38 percentage points of annualized stratum-month CPR is dominated by small-exposure strata with unstable realized rates, while the exposure-weighted RMSE --- the weighting the fit itself uses --- is 2.5 percentage points against an empirical mean CPR of 5.14\%, half the mean. Neither weighting supports out-of-sample predictive power at the stratum-month level (and the ridge grid's two candidate fits are numerically identical on this holdout, so no selection hinges on the weighting). Second, the recovery is concentrated where the model is unsupervised: of the \$915.0 billion, \$331.0 billion accrues in the 19 QT months overlapping the training window (June 2022--December 2023), against \$337.5 billion empirical --- 98\% --- while \$584.1 billion accrues in the 23 post-training simulation months from January 2024 onward (21 of which have performance data; the empirical comparison uses those), against \$427.3 billion empirical --- 137\%; components are independently rounded, and the overshoot lives almost entirely out of sample. Third, the raw CPR correlation against the empirical path is $r = -0.444$ at zero lag --- more negative even than the ABM's $-0.318$ --- and \eqref{eq:pathA} omits calendar-month seasonality, one of the four factors of \citepos{richard1989} canonical prepayment function and a natural first suspect for a monthly path fit this poor. Re-estimating with eleven calendar-month dummies confirms strong, correctly signed seasonality (spring and early-fall log-effects up to $+0.43$ against a January base) and improves the simulated path's contemporaneous correlation from $-0.444$ to $-0.378$, with the peak moving from lag $-3$ to $-2$ and the recovery to 121.5\% --- real, but far from sufficient: the named first suspect accounts for only a modest share of the negative path fit. The production specification is unchanged this round --- adopting the month effects would break comparability with the frozen-run ledger mid-review --- and the seasonal specification is slated to become production at the next full re-freeze; until then it is reported alongside as robustness. A model that recovers 119.7\% of the dollar benchmark while mismatching the monthly path this badly may be recovering the aggregate partly through offsetting errors: over-predicting prepayment in some cohort-months and under-predicting in others in a way that nets out near the benchmark without the underlying mechanism being correctly timed. We therefore treat Path A's recovery as evidence about the aggregate level only. No timing claim rests on Path B either: its apparent levels advantage does not survive detrending (Section~\ref{sec:pathb}), and every claim this paper makes for Path B is a level-accuracy claim.

\subsection{Path B: Literature-Calibrated Microsimulation}\label{sec:pathb}

The second hazard path addresses competing risks, literature-calibrated elasticities, and a cleaner isolation of the Danish counterfactual, without reintroducing the utility-maximizing machinery of the ABM. It simulates a stratified sample of 75,000 loans (40,234 of which are still active when the QT window opens; Appendix~\ref{app:params}), each carrying static covariates (FICO, region, original LTV, coupon, and a stratum identifier matching the Path A specification) and dynamic state (loan age, balance, and a regime-specific rate gap). Each loan's monthly hazard of prepayment is specified as
\begin{equation}\label{eq:pathB}
h^{\mathrm{prep}}_{i,t} = \max\!\left\{\, \underline{h},\; h_0\!\left(a_{i,t}\right) \exp\!\left( \beta_1\, g_{i,t} + \beta_F\, z^{\mathrm{FICO}}_{i} + \beta_L\, z^{\mathrm{LTV}}_{i} + \beta_b\, B_{s(i),t} \right) \right\},
\end{equation}
with a competing hazard of default
\begin{equation}\label{eq:default}
h^{\mathrm{def}}_{i,t} = h^{\mathrm{def}}_0 \exp\!\big( \gamma_r\, \sigma_{i,t} + \gamma_F\, z^{\mathrm{FICO}}_{i} + \gamma_L\, z^{\mathrm{LTV}}_{i} \big), \qquad \sigma_{i,t} = \max\{0,\ -\tilde g_{i,t}\},
\end{equation}
where $\sigma_{i,t}$ is rate stress --- $\tilde g_{i,t} = g_{i,t}/100$ is the rate gap of \eqref{eq:pathB} expressed in decimal units, so $\sigma$ is the decimal magnitude of an out-of-the-money gap --- and $h^{\mathrm{def}}_0 = 3\times10^{-4}$ monthly. Every Path B parameter is fixed a priori in the configuration and none is estimated on the benchmark; Appendix~\ref{app:params} collects the full parameter vector with provenance, together with the delinquency transition matrix. Here $h_0(a_{i,t})$ is the baseline hazard at loan age $a_{i,t}$; $g_{i,t}$ is the rate gap --- the loan's coupon minus the prevailing market rate, in 100-basis-point units, negative when the borrower is locked in, so that the positive $\beta_1$ of \eqref{eq:beta1} suppresses prepayment exactly when the par-payoff penalty binds; $z^{\mathrm{FICO}}_{i}$ and $z^{\mathrm{LTV}}_{i}$ are standardized static covariates; $B_{s(i),t} \in [0,1]$ is the cohort-level burnout state of loan $i$'s stratum $s(i)$ --- the stratum's cumulative prepaid UPB as a share of its original UPB, capped at one, so $\beta_b = -0.5$ scales the hazard by $e^{-0.5 B} \in [0.61, 1]$; and $\underline{h}$ is the involuntary-turnover floor: the monthly single-month-mortality (SMM) equivalent of 4\% annual CPR, $\underline{h} = 1-(1-0.04)^{1/12} \approx 0.0034$, representing death, divorce, and forced relocation. The floor's level is anchored to the same empirical observation as the ABM's mobility calibration (Section~\ref{sec:method-frictions}): turnover on deeply out-of-the-money discount cohorts during 2023--2024, which lies \emph{inside} the June 2022--November 2025 evaluation window --- so while no Path B parameter touches the benchmark dollar figure, the specification's dominant level-setter is measured from in-window prepayment behavior, and the verification reading below must be taken net of it: the clean verification content is the +9.2-point marginal over the $\beta_1 = 0$ null, not the aggregate level itself. The floor binds in 36\% of the 1,683,124 evaluated loan-months in the production U.S. run, so a substantial share of the aggregate level rides on the floor rather than the elasticity --- consistent with the no-lock-in null of Section~\ref{sec:hazard-interp}, which shows the mechanical model recovers 97.8\% of the benchmark. Proportionality in \eqref{eq:pathB} accordingly holds above the floor, not through it. The baseline hazard follows the 100~PSA seasoning ramp --- CPR rising 0.2 percentage points per month of loan age to 6\% at month 30 --- applied at monthly frequency via $\mathrm{SMM} = 1-(1-\mathrm{CPR})^{1/12}$. The rate-gap coefficient is calibrated from \citet{liebersohn2024}, who estimate that a 100bp increase in the prevailing rate above the locked-in coupon reduces quarterly mobility probability by 5.5\% to 7.7\% depending on specification; we adopt the 6.5\% midpoint of that band as the ex-ante central scenario, with the band edges as sensitivity bounds\footnote{The central 6.5\% is the midpoint of that band, adopted as the ex-ante central scenario rather than a published point estimate, with 5.5\% and 7.7\% as sensitivity bounds. Quarterly probability shocks are mapped to the proportional-hazard coefficient $\beta_1$ via the survival-function transform in \eqref{eq:beta1}, evaluated at a baseline quarterly mobility probability of 6\% --- a turnover-proxy calibration auxiliary that affects only the conversion level, not the elasticity band itself, and demonstrably weakly: in the small-$P_q$ limit the transform reduces to $\beta_1 = -\ln(1-\delta) = 0.0672$ at the central $\delta$, and over $P_q \in (0, 0.12]$ --- from arbitrarily low turnover to double the adopted proxy --- the central $\beta_1$ ranges only from $0.0672$ to $0.0701$, so even a large error in $P_q$ moves the conversion by a few percent. \citet{fonseca2024} report a related but distinct estimand --- a roughly 9\% reduction in moving rates per percentage point of mortgage-rate delta --- which we cite as independent corroboration of the lock-in magnitude, not as a source of the band.}, converted to a monthly hazard via the survival-function transform
\begin{equation}\label{eq:beta1}
\beta_1 = -\ln\!\left[\frac{1-\big(1-(1-\delta)\,P_q\big)^{1/3}}{1-\big(1-P_q\big)^{1/3}}\right],
\end{equation}
where $\delta$ is the proportional quarterly mobility decline per 100bp of rate gap, $P_q = 0.06$ the baseline quarterly mobility probability, and the inner terms are the monthly single-month mortalities implied by constant-hazard compounding within the quarter, $1 - P_q = (1 - h_m)^3$; the central $\delta = 0.065$ gives $\beta_1 \approx 0.069$ (unrounded $0.0686$; an earlier draft printed the continuous-hazard form of this transform, which gives $0.0693$ --- both round to $0.069$, and Appendix~\ref{app:ledger} records the correction). The transform yields the dimensionless log-hazard-ratio implied by a quarterly probability decline --- it replaces a naive divide-by-three approximation, not a hazard itself --- and \eqref{eq:pathB} then applies $\beta_1$ linearly per 100bp of gap. At QT-era gaps of roughly $-380$bp this compounds the local quarterly estimate to a hazard ratio of about $0.77$ ($\approx e^{3.8 \times -0.069}$), a log-linear extrapolation beyond the range in which the source elasticity was estimated; we flag it as a functional-form assumption and test it directly: a piecewise-concave transform (full marginal effect over the first 200bp of lock-in, half beyond) moves the recovery from 107.0\% to 105.7\% (\$808.4 billion) and slightly improves the contemporaneous path correlation ($+0.281$ versus $+0.190$, peak lag unchanged at $-3$), so the log-linear extrapolation is not load-bearing. Each month, both hazards are evaluated for every active loan. In the production configuration prepayment is fractional: each active loan's balance is drained deterministically by $h^{\mathrm{prep}}$ --- correcting a description in earlier drafts that attributed prepayment to a per-loan uniform draw, which is a non-production discrete mode --- while a uniform draw governs the default hazard, sending hits into the delinquency pipeline (current loans to D30, already-delinquent loans to default), where they transition through the Markov delinquency-resolution matrix of Appendix~\ref{app:params}. A proportional rescaling guards the case $h^{\mathrm{prep}} + h^{\mathrm{def}} > 1$; it is a safety clamp, not a competing-risks construction, and it never binds in production --- across all 1,683,124 evaluated loan-months per regime the maximum hazard sum is 0.016, two orders of magnitude below the threshold, at which level the exact exponential cause-allocation construction $1-\exp(-(\lambda_p+\lambda_d))$ with shares $\lambda_p/(\lambda_p+\lambda_d)$ differs from ours only at second order in the hazards. Prepaid balances settle to simulated SOMA cash flow in the month they occur; no settlement delay is modeled on this path.

Aggregation from the 75,000-loan sample to the SOMA book applies balance weights: simulated prepaid UPB is scaled by the ratio of the Fed's MBS book to the sample balance, with no reweighting toward Ginnie Mae collateral (20.4\% of book face) or vintages outside 2017--2021 --- the Freddie-estimated hazards are extrapolated to those segments as-is. Two bounds discipline this extrapolation. Along the coupon dimension, Section~\ref{sec:robustness-hybrid}'s full-book reweighting (sample WAC 3.9\% mapped to the book's 2.49\%) moves Path B from 107.0\% to 109.1\%, a composition effect of about two percentage points. Along the agency dimension, Ginnie ``prepayments'' include servicer buyouts of seriously delinquent loans, a mechanically different process from Freddie voluntary speeds, and an external bound now disciplines the extrapolation: Ginnie Mae's fixed-rate speeds ran 1.2--2.8 percentage points above Freddie Mac's across the window (5.2\% versus 4.0\% in February 2023; 10.3\% versus 7.5\% in May 2025; Ginnie Mae Global Markets Analysis Report, Recursion data), with the involuntary buyout channel (CDR 2.1\% versus 0.4\% in May 2025) supplying most of the structural difference. Treating the 20.4\% Ginnie face share as conventional therefore mis-states book-wide CPR by at most $0.204 \times (1.2\text{--}2.8) \approx 0.24$--$0.57$ percentage points; at the null-to-central sensitivity of roughly \$83 billion per CPR point, the composition error is bounded at \$20--\$47 billion (2.6--6.2\% of benchmark), signed toward overstating trapped liquidity. Burnout in this specification is deliberately cohort-level rather than per-loan: because an individual loan's prepayment is a one-time, absorbing transition rather than a continuous rate, a per-loan fractional burnout measure is not well-defined --- a category error the closed-population survivor-selection failure of Section~\ref{sec:abm-burnout} illustrates concretely. Instead, each stratum accumulates its cumulative prepaid UPB as a share of its original UPB (the $B \in [0,1]$ state defined under \eqref{eq:pathB}) after all prepayment transitions in a given month, broadcast to all of that stratum's remaining active loans before the next month's hazard is evaluated, and entering \eqref{eq:pathB} with the calibrated prior coefficient $\beta_b = -0.5$ (Appendix~\ref{app:params}); Path A's separately estimated burnout coefficient ($-0.13$, statistically zero) plays no role here. Because the simulated burnout state is endogenous, we verify the channel is load-bearing rather than decorative with a $\beta_b = 0$ rerun on the production seed: trapped liquidity moves by \$6.6 billion ($-0.9$pp of benchmark, to 106.2\%; endpoint rounding makes this $-0.8$pp if computed from the rounded shares), with the peak lag unchanged at $-3$. The FICO and LTV covariates receive the same treatment: a $\beta_F = \beta_L = 0$ rerun moves the recovery to 108.3\% ($+1.3$pp), and a rerun under coefficients estimated Path-A-style from the panel moves it to 101.3\% ($-5.7$pp); Appendix~\ref{app:params} reports the estimation, whose imprecise point estimates are consistent with the priors' turnover-regime signs without rejecting alternatives. Production retains the ex-ante priors --- the estimates are too imprecise to license replacing a pre-specified calibration mid-stream --- and we flag the covariate block as the specification's largest single sensitivity: the two endpoints correspond to 99.2\% and 92.2\% on the shared accounting basis (the common netting makes the conversion exact), a swing comparable in size to the lock-in marginal itself.

Two Danish counterfactual mechanisms exist in this framework's pipeline, and only one is live in production. The first is a \emph{rate-gap swap}: holding the baseline hazard, rate-gap coefficient, covariates, and random draw logic fixed, the Danish leg expresses the market-value buyback in rate-equivalent units, so that the buyback discount properly offsets the locked-in spread --- $g_{\mathrm{DK}} = g_{\mathrm{US}} + \Delta y_{\mathrm{buyback}} \approx 0$ whenever the loan is out of the money, and $g_{\mathrm{DK}} = g_{\mathrm{US}}$ when rates have fallen and the buyback price is capped at par (at the representative 3.0\% coupon / 6.8\% market state the U.S. gap is $-380$bp while the Danish effective gap is zero, so the U.S. hazard is suppressed by the full locked-in spread while the Danish hazard resets toward its baseline seasoning curve). This swap defines the \emph{mechanism-substitution} variant --- the specification behind the superseded 47.1\% Danish CPR of Section~\ref{sec:abm-danish} --- and it is \emph{not} the production mechanism. The production Danish leg is the Berger-recalibrated specification adopted in Section~\ref{sec:abm-danish}: it replaces the rate-gap swap with the imported flat moving elasticity, so the production Danish hazard does not respond to the buyback discount at all. The swap machinery remains in the codebase and its PV-versus-par classification is still computed each month, but no production quantity consumes its output --- which is why the discount-rate diagnostic of Section~\ref{sec:method-abm} moves the production Danish leg by exactly \$0.00 billion.

Simulated on the 40,234 window-start survivors of the 75,000-loan stratified sample drawn from the same Freddie Mac 2017--2021 origination universe (Appendix~\ref{app:params} reports the full attrition accounting), this specification recovers \$818.5 billion, or 107.0\%, of the empirical benchmark, with a raw CPR correlation against the empirical path of $r = +0.190$ at zero lag and a peak cross-correlation of $r = +0.404$ at lag $-3$. Throughout, the cross-correlation at lag $k$ pairs the simulated series at month $t$ with the empirical series at month $t+k$ --- equivalently, the empirical series at month $t$ with the simulated series at month $t-k$ (both series demeaned over the overlap; boundary months dropped) --- so at the peak $k=-3$ the simulated value at month $t$ is matched with the empirical value from three months earlier. An earlier draft printed this definition with the two series interchanged, which would have implied the opposite direction; the direction claims themselves were always computed under the implemented convention (Appendix~\ref{app:ledger}). A direction check with synthetic data confirms that under this implementation a peak at $k=-3$ means the simulated path \emph{trails} the empirical path by three months --- the empirical series moves first. At this peak, the empirical SOMA-derived CPR path moves approximately three months ahead of the simulated path (a peak selected across the seven examined lags, so its nominal significance overstates; no inference rests on this single coefficient). An earlier specification applied the rate-gap coefficient $\beta_1$ to the decimal rate gap rather than to the gap measured in 100-basis-point units in which $\beta_1$ is expressed, which left the calibrated lock-in elasticity numerically inert (a multiplier of approximately $\times$1.0026 rather than the intended $\times$0.77 at a typical QT-era rate gap); under that earlier, mis-specified version, the elasticity channel contributed nothing, and the \$747 billion figure it produced reflected only the PSA baseline, involuntary floor, and burnout terms. Correcting this unit mismatch is what produces the level and band reported here (see Section~\ref{sec:robustness-pipeline}).

Two things follow from this calibration choice. First, unlike Path A, whose rate-gap and burnout coefficients are estimated directly from the Freddie Mac sample, Path B's rate-gap elasticity is imported unmodified from \citet{liebersohn2024} and never touches the \$764.7 billion benchmark during calibration. Path B is accordingly not a second attempt to fit our benchmark --- it is a structural verification exercise: it asks whether an elasticity estimated on an entirely different empirical setting, dropped into a PSA-baseline hazard with competing risks and no tuning against our target, reproduces the Federal Reserve's actual QT runoff shortfall. That it does so within 8 percentage points of the benchmark is evidence bearing on the external validity of the Liebersohn--Rothstein elasticity itself --- conditional on the marginal contribution's seed-level uncertainty reported in Section~\ref{sec:hazard-interp}, which is small --- not merely on Path B's specification --- but the timing is not part of that evidence: as Section~\ref{sec:hazard-interp} reports, a rerun with the elasticity set to zero reproduces the same three-month lead, so the lead cannot discriminate the elasticity from the mechanical model, and the elasticity's identified content is its marginal recovery over that null together with the monotone response across the calibration band. Second, propagating the full Liebersohn--Rothstein 5.5\%--7.7\% sensitivity band through the simulation yields a recovery range of \$809.9 billion to \$827.7 billion (105.9\%--108.2\% of benchmark under the standalone scorer; 96.8\%--99.1\% on the shared accounting basis of Section~\ref{sec:robustness-hybrid}), with the peak cross-correlation lag stable at $-3$ across the entire band (peak $r$ between $+0.371$ and $+0.423$; given the moving-block intervals above and the seven-lag search, these differences carry no evidential weight). Stronger assumed mobility suppression therefore maps monotonically to more trapped liquidity, as expected, and the central 107.0\% estimate should be read as a point within this now-quantified interval.

We correct here a directional claim made in earlier drafts: the peak alignment does \emph{not} show the simulated path leading SOMA cash flow. Under the implemented convention (verified against synthetic data), the empirical path moves first --- a settlement-delay mechanism, in which loan-level prepayment decisions show up in SOMA cash two to three months later, predicts the \emph{opposite} direction of offset. We therefore attach no settlement interpretation to the timing alignment.

A direct input-timing diagnosis rules out the dating of the rate input and narrows --- but does not isolate --- the mechanism. In monthly \emph{changes}, the simulated CPR responds to the contemporaneous rate change with the correct sign and near-mechanical strength ($\Delta$-correlation $-0.94$ at lag 0), while the empirical CPR's month-to-month changes show no detectable rate response ($+0.25$, inside the $\pm 0.31$ zero band at $n = 41$): at monthly frequency the empirical path's variation is dominated by non-rate factors --- calendar seasonality above all --- that the simulation does not model.

The same caution applies to the paths themselves. Path B's $+0.190$ lag-0 levels correlation is carried by the window's shared trend: in first differences it is $-0.23$ and after linear detrending $-0.20$. The same constructions for the other estimators give $-0.20$ and $-0.30$ for the ABM (from $-0.32$ in levels) and $-0.21$ and $-0.29$ for Path A (from $-0.44$ in levels) --- detrending collapses the three estimators' apparent differences into a cluster between $-0.20$ and $-0.30$. Against the i.i.d.\ $1.96/\sqrt{n}$ reference ($\pm 0.30$ at $n = 42$ for the detrended series; $\pm 0.31$ at $n = 41$ for first differences), Path A ($-0.29$) and Path B ($-0.20$) are inside under both empirical-series conventions, while the ABM's detrended value straddles the band's edge across them ($-0.301$ common-series, $-0.303$ against its own back-out); indistinguishability verdicts therefore rest on the wider moving-block bands of Table~\ref{tab:estimators}, under which none of the three is distinguishable from zero. Once the shared trend is removed, no estimator's contemporaneous co-movement is positive and none is distinguishable from zero; Path B's distinction is the accuracy of its aggregate level, not its timing. (In levels, which are trend-contaminated over the window and reported only for completeness, the simulated path co-moves most strongly with the rate from four months earlier while the empirical path is contemporaneous.) No single channel carries the offset: the $\beta_b = 0$ and $\beta_1 = 0$ ablations both leave the peak at $-3$, and re-running the simulation with the rate input shifted by one to three months in either direction never moves the peak to lag 0 (it reaches $-1$ under a two-to-three-month input lead) while moving recovered trapped liquidity by less than \$1.5 billion across the entire scan. The offset is therefore a property of the fitted path's smooth aggregate dynamics set against a seasonal empirical series --- mechanism not isolated --- and Path A's seasonal re-estimation moving its own peak from $-3$ to $-2$ (Section~\ref{sec:patha}) is consistent with seasonality carrying part of it.

A moving-block bootstrap (block length 6, 4,000 replications) puts the 95\% interval on the peak-lag correlation at $[-0.21, +0.49]$ and on the lag-0 correlation at $[-0.13, +0.44]$; with seven lags searched, individual band-point differences in peak $r$ carry no evidential weight. These intervals are not sensitive to the block-length choice in any verdict-relevant way: at block lengths 4 and 8 the peak-lag endpoints move to $[-0.27, +0.39]$ and $[-0.22, +0.51]$ and the lag-0 endpoints to $[-0.11, +0.45]$ and $[-0.17, +0.42]$, every interval still containing zero. The full specification, including loan sample construction, runs in approximately 15 seconds from cached inputs per band point, or approximately two minutes for a full rebuild across 75,000 loans, 42 months, and both institutional regimes.

\subsection{Interpretation}\label{sec:hazard-interp}

Table~\ref{tab:estimators} summarizes the three estimators evaluated in this paper against the common \$764.7 billion empirical benchmark.

\begin{table}[t]
\centering
\small
\begin{threeparttable}
\caption{Estimator comparison against the \$764.7 billion empirical trapped-liquidity benchmark.}
\label{tab:estimators}
\begin{tabular}{@{}>{\raggedright\arraybackslash}p{3.2cm}>{\raggedleft\arraybackslash}p{2.4cm}>{\raggedleft\arraybackslash}p{2.4cm}>{\raggedright\arraybackslash}p{2.3cm}>{\raggedright\arraybackslash}p{4.0cm}@{}}
\toprule
Estimator & Trapped liquidity & \% of benchmark & CPR path corr. (lag) & Path diagnostics \\
\midrule
Empirical benchmark (SOMA, phased cap, 30+15yr) & \$764.7B & 100.0\% & --- & mean CPR 5.14\% \\
Production ABM (Section~\ref{sec:abm-results}) & \$91.0B & 11.9\% & $r = -0.318$ (0) & mean CPR 11.68\%; raw $R^2$ = $-6.984$ \\
Hazard Path A --- stratum-month Poisson PML & \$915.0B & 119.7\% & $r = -0.444$ (0) & mean CPR 3.51\%; holdout RMSE $\approx$ 38pp \\
Hazard Path B --- literature microsimulation & \$818.5B & 107.0\% (band 105.9--108.2\%) & $r = +0.404$ ($-3$) & mean CPR 4.76\%; $r = +0.190$ at lag 0; 3-month offset (empirical leads; reference only, Section~\ref{sec:pathb}) \\
Hazard Path B --- concave-gap variant (Section~\ref{sec:pathb}) & \$808.4B & 105.7\% & $r = +0.281$ (0) & piecewise transform, kink 200bp; robustness on the log-linear extrapolation \\
Cross-design ABM, real covariates, recalibrated (Section~\ref{sec:robustness-crossdesign}) & \$453.5B & 59.3\% & $r = -0.336$ (0) & mean CPR 8.14\%; no positive alignment at any lag \\
Cross-design ABM, real covariates, frozen calibration (Section~\ref{sec:robustness-crossdesign}) & \$159.5B & 20.9\% & $r = -0.311$ (0) & mean CPR 11.63\%; no positive alignment at any lag \\
\bottomrule
\end{tabular}
\begin{tablenotes}[flushleft]\footnotesize
\item The table reports the three primary estimators and the two Section~\ref{sec:robustness-crossdesign} cross-design variants; the Danish counterfactual specifications are collected in Table~\ref{tab:danish}. Dashes denote quantities that are not applicable to the row. Peak-lag entries select the maximum-magnitude cross-correlation (max $|r|$) within a $\pm$3-month window; under a most-positive-$r$ convention the production ABM's best alignment is $+0.19$ at lag $-3$ (Section~\ref{sec:abm-results}) --- the two conventions summarize the same lag profile, and Path B's lag $-3$ peak satisfies both.
\item Path B is the only estimator whose aggregate recovery is accurate on the benchmark-consistent basis (97.9\%; 100.0\% composed). Its raw lag-0 correlation ($+0.190$) is carried by the window's shared trend: in first differences ($-0.23$) and after linear detrending ($-0.20$) the simulated--empirical co-movement is indistinguishable from zero ($\pm 0.31$ at $n = 41$), as it is for every estimator. The two cross-design rows substitute real Freddie Mac structural covariates into the ABM under two calibration choices; neither reproduces the empirical path's timing (no estimator does; see above), but the recalibrated variant's dollar recovery exceeds our own ex-ante threshold for undercutting the paradigm interpretation (Sections~\ref{sec:robustness-crossdesign} and~\ref{sec:conclusion}). ABM figures cite frozen run \texttt{run-2026-07-04-15yr-foldin} (footnote~\ref{fn:manifest}). Moving-block bootstrap 95\% intervals (block length 6): ABM lag-0 $[-0.60, +0.04]$; Path A lag-0 $[-0.68, -0.02]$; Path B lag-0 $[-0.13, +0.44]$, peak lag $[-0.21, +0.49]$ --- no positive lag-0 coefficient is distinguishable from zero --- Path A's contemporaneous levels correlation is \emph{negative} on the levels construction (interval excluding zero at block length 4; marginal at 6 --- upper endpoint $-0.02$ at the production seed, $-0.00$ under the committed tool's reconstruction --- and at 8: $[-0.66, +0.04]$), which strengthens the offsetting-errors reading of its aggregate fit (detrended, all three estimators cluster at $-0.20$ to $-0.30$; Section~\ref{sec:pathb}) --- and no timing evidence is claimed for any estimator. Block-length sensitivity (4/6/8) changes no zero-inclusion verdict in this table; the Path A significance caveat just given is the only endpoint that crosses.
\item Correlation-column convention: the ABM row's correlations are computed against the ABM pipeline's own empirical back-out (cohort-weighted, term-aware scheduled amortization), while the hazard rows use the hazard-side back-out; against the common hazard-side series the ABM lag-0 value is $-0.316$ --- a third-decimal, no-verdict difference. Section~\ref{sec:method-benchmark} notes the two constructions.
\item Shared-accounting-basis equivalents (Section~\ref{sec:robustness-hybrid}; the \$69.6B curtailment netting is common to all rows, so subtract 9.1pp): Path A 110.6\%, Path B 97.9\% (band 96.8--99.1\%), concave variant 96.6\%, composed Path B 100.0\%. Path A's point estimate carries no useful forward precision: the refit-and-resimulate 95\% percentile interval is $[-\$4{,}459.8, +\$1{,}190.1]$ billion (median \$922.6B, IQR \$900.1--\$1{,}130.6B; all standalone-basis --- subtract the common \$69.6B netting for shared-basis equivalents), itself a lower bound on propagated uncertainty since the age spline and stratum fixed effects are held at the production fit.
\end{tablenotes}
\end{threeparttable}
\end{table}

The two hazard paths, despite sharing no household-level decision machinery whatsoever with the ABM, together land far closer to the empirical benchmark --- 97.9\% and 110.6\% of it on the benchmark-consistent shared accounting basis of Section~\ref{sec:robustness-hybrid}, 107.0\% and 119.7\% under their standalone scorers --- than any synthetic-population ABM specification tested in this paper, including the falsified burnout attempt. Overshooting less, however, is not validation, and we do not present it as such: Path A overshoots the benchmark by nearly \$150 billion with a monthly path fit no better than the ABM's, and Path B's recovery, now reported across its full elasticity sensitivity band (\$809.9B--\$827.7B, 105.9\%--108.2\%; Section~\ref{sec:pathb}), reflects an elasticity imported unmodified from \citet{liebersohn2024} rather than one fit to this paper's own benchmark --- a structural verification of that external estimate, not a second attempt to hit our target. The comparative claim we defend is therefore narrower and more precise than ``the hazard framework fits'': it is that Path B --- a loan-level survival structure with no household-utility content --- lands within 2.1\% of the aggregate benchmark on the benchmark-consistent accounting basis under the ex-ante covariate priors (92.2\%--99.2\% across covariate calibrations, Section~\ref{sec:pathb}; within 0.05\% in the composed configuration under those priors, approximately 94.3--101.3\% composed across covariate calibrations), with Path A, whose forward simulation carries no useful propagated precision (Section~\ref{sec:patha}), agreeing on sign and order of magnitude at 110.6\% under a specification that shares data but no estimated parameters, while no synthetic-population household-choice specification we tested lands within 45\% of it. We attach no timing credential to this: Path B's raw $+0.190$ lag-0 correlation is carried by the window's shared trend --- in first differences and after detrending the co-movement is indistinguishable from zero (Section~\ref{sec:pathb}) --- and the moving-block intervals of Table~\ref{tab:estimators} counsel against weighing any single timing coefficient, including Path B's lag $-3$ peak, whose interval spans zero. Timing alone, however, does not identify the lock-in channel: a rerun of Path B with the lock-in elasticity set to zero also peaks at lag $-3$ while still recovering 97.8\% of the benchmark under the standalone scorer (88.7\% on the shared layer; the margin is basis-invariant because the accounting netting is common to both runs), so the elasticity's identified contribution is its marginal recovery --- +\$70.3 billion, or +9.2 percentage points, comparable in magnitude to the specification's covariate-prior sensitivity (Section~\ref{sec:pathb}), a margin whose seed-level uncertainty we now report directly: across 10 paired reruns with different draw seeds (central and null sharing each seed), the margin is \$70.3 billion and varies by less than \$0.01 billion, because the production engine drains prepayments fractionally and the draw seed touches only the small delinquency channel; the uncertainty that matters for the margin is the calibration band and the structural assumptions, not simulation noise --- and the monotone response across the Liebersohn--Rothstein band, not the lag structure itself. We qualify this claim to synthetic-population specifications deliberately: Section~\ref{sec:robustness-crossdesign} reports a real-covariate cross-design variant of the ABM that does land within 45\% of the benchmark (59.3\%), which is precisely the result that complicates rather than confirms the paradigm interpretation this paragraph would otherwise support.

Because the hazard framework's only behavioral content is a set of loan-level covariates, a cohort-level burnout term, and literature-calibrated elasticities --- with no analogue to mobility desire, loss aversion, or utility maximization --- its comparative success is evidence, though, given the unresolved data-source asymmetry of Section~\ref{sec:robustness-scale}, not proof, for the loan-level and institutional mechanisms (heterogeneity in FICO, LTV, and geography; servicer forbearance and pipeline delay; and cohort-level burnout) that the synthetic-population ABM of Section~\ref{sec:abm} does not represent. The scale-sensitivity replication in Section~\ref{sec:robustness-scale} rules out the ABM/hazard sample-size difference as the driver of this gap; the data-source difference is addressed only partially by the cross-design test in Section~\ref{sec:robustness-crossdesign}, which finds that endowing the ABM with real structural covariates changes its explanatory share substantially --- enough, under one calibration choice, to raise the question of whether ``cannot represent'' is the right description at all. The rate-gap coefficient is stable in sign across both hazard specifications, indicating that the underlying par-payoff penalty is doing real work; household rationality is not irrelevant, but the aggregation structure of loan-level survival is what a model of this shortfall must get right at minimum, and Section~\ref{sec:robustness-crossdesign} leaves open how much more it must get right beyond that. This sharpens the trade-off named in this paper's title, with the qualification Section~\ref{sec:robustness-crossdesign} adds: the price of the TBA market's homogeneity appears to be paid at least in part through loan-level and institutional frictions --- FICO- and geography-driven heterogeneity, servicer pipeline delay, cohort-level burnout --- that the synthetic-population household model does not represent and that a market-value repurchase design, as Section~\ref{sec:abm-danish}'s counterfactual suggests, would substantially relieve.

\section{Discussion}\label{sec:discussion}

Sections~\ref{sec:abm} and~\ref{sec:hazard} established the paper's central contrast: household simulation, rational or behaviorally extended, does not close the gap to the empirical benchmark, while a structurally independent hazard framework --- discarding household utility altogether, on a substantially larger and non-synthetic sample --- recovers 97.9--110.6\% of it on the benchmark-consistent accounting basis, its microsimulation form landing within 2.1\% --- with the caveat that no estimator's monthly path co-movement, Path B's included, survives removal of the window's shared trend (Section~\ref{sec:pathb}). Taken together, and bearing in mind the data-source caveat above, we view three candidate mechanisms as plausible contributors to that gap. First, loan-level heterogeneity in credit score, geography, and life circumstance means that mobility decisions are not a smooth function of coupon and market rate alone; the hazard framework's 295 stratum fixed effects, which the ABM's coupon-cohort buckets cannot replicate, are one mechanism that lets Path A track pool-level heterogeneity the ABM does not represent. Second, servicer behavior --- loss mitigation, forbearance, and loan modification practices that vary by servicer and were unusually active during the higher-rate, higher-delinquency-risk period of the QT window --- operates on prepayment and default outcomes through channels outside a borrower-side decision model; Path B's explicit competing-risks structure and Markov delinquency pipeline are a direct, if partial, representation of this channel. Third, prepayment burnout is representable at the cohort level (as in both hazard paths) but not at the level of a closed, per-agent household population, exactly as Section~\ref{sec:abm-burnout}'s falsified survivor-selection experiment demonstrated.

\subsection{The Dual Systemic Fallout}\label{sec:discussion-fallout}

While the macroeconomic literature frequently focuses on the housing mobility implications of the lock-in effect, its transmission into the fixed-income sector creates a dual-pronged systemic risk.

\paragraph{Central bank impact.} The extension of MBS durations starves the SOMA of its scheduled cash flows, breaking the predictability required for effective Quantitative Tightening; Section~\ref{sec:robustness-wal} expresses this extension in approximate weighted-average-life units.

\paragraph{Private bank fragility.} Unlike the Fed, commercial banks must mark their portfolios to market. The structural duration extension of these low-coupon MBS assets generated hundreds of billions in unrecognized losses across the banking sector. \citet{jiang2023} show that duration extension alone was not sufficient to precipitate failure; rather, it was the interaction of mark-to-market losses with a high share of uninsured deposits that made specific regional institutions vulnerable to self-fulfilling runs, as uninsured depositors withdrew funds upon recognizing that reported asset values no longer covered liabilities. This distinction matters for our argument: the lock-in effect generates the duration extension that is a necessary precondition for this fragility channel, even though it is not, on its own, a sufficient condition for bank failure.

\paragraph{Labor market misallocation.} Beyond its transmission into fixed-income markets, the lock-in effect imposes a real economic cost by preventing households from relocating toward higher-productivity labor markets. \citet{hsieh2019} estimate that housing supply constraints in high-productivity U.S. cities impose substantial aggregate output losses through this kind of spatial misallocation of labor; their headline magnitudes are cumulative level effects over 1964--2009 and were revised in a published correction, so we cite the mechanism rather than a point estimate. While their analysis concerns supply-side constraints on housing availability, the demand-side mobility constraint we quantify here operates through the same channel: households who cannot economically justify relocating cannot respond to geographic differences in labor demand, regardless of housing supply elsewhere. The \$764.7 billion empirical shortfall this paper documents therefore likely understates the total welfare cost of the par-payoff rule, since it captures only the Federal Reserve balance-sheet channel and does not include forgone gains from labor reallocation.

\subsection{Evaluating Alternative Securitization Policies}\label{sec:discussion-policy}

In response to the mobility collapse, U.S. policymakers, including the Federal Housing Finance Agency (FHFA), proposed expanding assumable and portable mortgages to alleviate the lock-in effect. Fixed-income analysis suggests these proposals address only part of the problem: while portability effectively addresses consumer mobility constraints, it explicitly prevents the loan from prepaying upon relocation. The same holds for assumability, with one added wrinkle: an assumed loan also fails to prepay --- the below-market coupon transfers to the buyer and the extension stays on the investor's book --- and the buyer must separately finance the gap between the purchase price and the assumed balance, typically at prevailing rates, which limits take-up. Both instruments therefore relocate rather than eliminate duration extension risk --- solving the microeconomic mobility constraint while cementing the macroeconomic extension risk permanently onto the investor's balance sheet.

A second lever, which we flag explicitly as speculative design space rather than as a finding this paper's estimates support, operates on the policy instrument rather than the mortgage contract: the design of the redemption cap itself. The shortfall we measure is defined against a fixed \$35 billion monthly cap that locked-in prepayment behavior predictably cannot reach, so a cap indexed to the observable state of the outstanding stock --- for example, to the share of the book sitting deeply below market coupon --- or a standing policy of substituting Treasury runoff or outright MBS sales when passive runoff undershoots, would convert extension risk from a silent shortfall into an explicit, priced policy choice. Nothing in our estimates identifies the welfare or market effects of such designs; active sales in particular would realize the mark-to-market losses that passive runoff defers. We record the option to mark where the design space lies, not to recommend a point in it.

\section{Robustness Checks}\label{sec:robustness}

This section collects the methodological checks referenced throughout the paper. With one exception, none of them changes the paper's central result --- that synthetic-population household choice explains a small share of the Federal Reserve's trapped liquidity while the loan-level hazard framework overshoots it at the aggregate level. The exception is stated where it occurs: the recalibrated cross-design variant of Section~\ref{sec:robustness-crossdesign} exceeds our ex-ante threshold and qualifies the household-choice reading at the aggregate level. We report these checks here rather than alongside their first mention so that the main narrative is not repeatedly interrupted to pre-empt objections.

\subsection{Scale Sensitivity}\label{sec:robustness-scale}

The ABM's 10,000-household population and the hazard framework's 75,000-loan sample differ in both scale and data source: the ABM draws a synthetic population calibrated to aggregate FRED medians, while the hazard framework is estimated directly on real Freddie Mac loan-level records. To test whether population size alone drives the gap between the two frameworks, we re-ran the ABM's core household mechanic --- log-normal mobility-desire heterogeneity mapped through a rate-gap-driven CPR function --- at N = 10,000 and N = 75,000 across 50 seeds each. The mean simulated CPR is 15.616\% at both population sizes (difference 0.0003 percentage points), while the 95\% confidence interval narrows by a factor of 2.4$\times$ at the larger $N$ --- of the same order as, though about 12\% short of, the 2.74$\times$ an idealized $\sqrt{N}$ law would predict ($\sqrt{7.5} \approx 2.74$), consistent with a non-i.i.d.\ component of seed noise; we do not lean on the exact ratio. Population size changes the precision of the ABM's estimate, not its central tendency. The gap between the ABM's 11.9\% and the hazard framework's 107.0--119.7\% of the empirical benchmark is therefore not attributable to the difference in sample size between the two estimators. It does not follow, however, that the gap is therefore attributable to modeling paradigm alone: that distinction cannot be closed by resampling, and Section~\ref{sec:robustness-crossdesign} takes it up directly with the completed cross-design pair.

\subsection{Benchmark Sensitivity to Data Source and Cap-Schedule Assumptions}\label{sec:robustness-benchmark}

The \$764.7 billion empirical benchmark (Section~\ref{sec:method-benchmark}) is constructed from the New York Fed's SOMA current-face-value series under the actual phased redemption cap schedule, netting realized roll-off against the cap within the active QT window. We treat net accumulation, rather than one-sided accumulation of shortfall months only, as the correct construction because it is the only version comparable to the Danish counterfactual of Section~\ref{sec:abm-danish}, where simulated roll-off routinely exceeds the cap; a one-sided construction would overstate the empirical benchmark relative to that counterfactual and bias the recovered percentages reported for both estimators. Researchers extending this work with alternative SOMA vintages or a different cap-schedule assumption should expect the benchmark level, and therefore the recovered percentages in Sections~\ref{sec:abm} and~\ref{sec:hazard}, to shift accordingly; the qualitative ranking of the three estimators reported in Table~\ref{tab:estimators} is not sensitive to this choice, since the ABM's shortfall against any plausible benchmark construction remains an order of magnitude smaller than either hazard path's. Section~\ref{sec:method-benchmark} separately addresses the conceptual status of the cap-relative construction itself. A related window-construction pitfall is worth recording as a reproducibility note: an earlier draft filter that accumulated from the QT start date without an upper bound inadvertently included eight post-QT months (December 2025 onward), drifting the headline CPR means by roughly 0.3--0.5 percentage points while leaving the dollar totals unchanged, since the dollar aggregation already applied the bounded active-QT mask. That asymmetry --- path statistics silently sensitive to a window filter that dollar aggregates are robust to --- is why every headline figure in this paper cites a single tagged run manifest (footnote~\ref{fn:manifest}).

\subsection{Behavioral Extension Mechanics}\label{sec:robustness-extensions}

This subsection reports the mechanism behind each of the three behavioral extensions summarized in Section~\ref{sec:abm-results}. As Section~\ref{sec:method-frictions} describes, loss aversion steepens the penalty every household faces, which on its own would suppress simulated mobility below the historically observed 4--5\% involuntary-turnover floor; to preserve that floor, we re-calibrate the mobility-desire parameter upward, nearly tripling it. That recalibration mechanically inflates the model's overall responsiveness to rate gaps, so the jump to 54.9\% is an artifact of holding the turnover floor fixed, not a validation of loss aversion as a mechanism. The second extension, income-scaled curtailment, pulls the estimate down to 45.1\%; because curtailment is a strictly additive drag on prepayment, this decline is exactly what we predicted ex ante from a diagnostic correlation between the ABM's over-prediction and disposable-income growth (Section~\ref{sec:method-frictions}), and is not a surprise. The third extension, multi-vintage coupon composition, pulls the estimate down further, to 33.7\% --- the opposite of what we expected. A lower weighted-average coupon widens the gap to prevailing market rates, and our hypothesis was that this should increase predicted trapped liquidity, not reduce it. That the result runs counter to the hypothesis is itself informative: it suggests the ABM's CPR mechanism is not translating a wider coupon gap into slower prepayment the way the household-choice framework predicts it should. We cannot rule out that a differently specified household model would perform better, but that possibility does not bear on the results reported in Section~\ref{sec:abm}.

\subsection{Cross-Design Test: Real Covariates in the ABM}\label{sec:robustness-crossdesign}

Section~\ref{sec:method-estimators} flags an unresolved confound between this paper's two estimators: they differ in modeling paradigm (household choice versus loan-level survival) and in data source (synthetic population versus real Freddie Mac loan records) simultaneously, so a low ABM share and a high hazard-framework share are equally consistent with either explanation. To make partial progress on separating them, we re-ran the ABM's decision rule with structural covariates --- coupon, loan age, and original LTV --- drawn directly from the same 75,000-loan Freddie Mac sample used throughout Section~\ref{sec:hazard}, in place of the log-normal, FRED-median-calibrated draws used elsewhere in this paper. Behavioral draws (mobility desire, loss aversion, patience, transaction cost) remain synthetic throughout, since the Freddie Mac data contains no analogue for these constructs; this test therefore isolates the effect of real structural heterogeneity, not a fully real population, and should be read with that scope in mind.

Because the ABM's mobility-scale parameter is calibrated to hold a real-world involuntary-turnover floor fixed (Section~\ref{sec:method-abm}), swapping in a population with different LTV and coupon moments raises a calibration question that a naive covariate swap does not resolve on its own: should the turnover-floor anchor be re-derived against the new population's moments, or held fixed at its original value? We report both variants rather than choosing one. The recalibrated variant applies the recalibration rule of Section~\ref{sec:method-frictions}, re-anchoring the mobility-scale parameter to reproduce the same 4--5\% turnover floor under the real population's moments; the frozen variant suspends that rule and holds the original synthetic-population calibration unchanged.

Before running either variant, we specified ex ante the following interpretive thresholds: a recovery above 50\% of the empirical benchmark would undercut the paradigm interpretation advanced in Sections~\ref{sec:abm}--\ref{sec:conclusion}; a recovery below 35\% would corroborate it, consistent with the synthetic-population ABM's own specifications; and a recovery between 35\% and 50\% would be ambiguous.

The recalibrated variant recovers \$453.5 billion, or 59.3\%, of the benchmark --- above our ex-ante 50\% threshold. The frozen variant recovers \$159.5 billion, or 20.9\%, within the corroborating band. Table~\ref{tab:crossdesign} summarizes both against the synthetic-population control.

\begin{table}[t]
\centering
\small
\begin{threeparttable}
\caption{Cross-design test results: real Freddie Mac structural covariates substituted into the ABM under two calibration choices, against the synthetic-population production ABM.}
\label{tab:crossdesign}
\begin{tabular}{@{}p{4.0cm}lll@{}}
\toprule
 & Recalibrated (primary) & Frozen (robustness) & Synthetic control \\
\midrule
Mobility-scale parameter & 36,086 & 43,883 & 43,883 \\
Trapped liquidity & \$453.5B (59.3\%) & \$159.5B (20.9\%) & \$91.0B (11.9\%) \\
Mean CPR & 8.14\% & 11.63\% & 11.68\% \\
CPR path corr. (lag 0) & $-0.336$ & $-0.311$ & $-0.318$ \\
Peak lag & 0 (no positive alignment) & 0 (no positive alignment) & 0 \\
\bottomrule
\end{tabular}
\begin{tablenotes}[flushleft]\footnotesize
\item Peak lag follows the max-$|r|$ convention of Table~\ref{tab:estimators}. Under the alternative most-positive-$r$ convention, the synthetic control's best alignment is $+0.195$ at lag $-3$ --- the ``$+0.19$ at the three-month lead'' of Section~\ref{sec:abm-results}; the two conventions summarize the same cross-correlation profile, and the substantive conclusion (no positive contemporaneous alignment --- which, after Section~\ref{sec:pathb}'s detrending result, no estimator including Path B shows) holds under either.
\end{tablenotes}
\end{threeparttable}
\end{table}

Two results stand out. First, by our own ex-ante criterion, the primary recalibrated variant's 59.3\% recovery moves against the paradigm interpretation: with real structural covariates and a like-for-like turnover anchor, the household-choice ABM explains far more of the benchmark than the synthetic-population specifications throughout Sections~\ref{sec:abm} and~\ref{sec:robustness-scale} suggested. Second, the \$294 billion gap between the recalibrated and frozen variants --- produced by nothing other than how the turnover-floor anchor is set, with the underlying covariates held fixed --- indicates that the ABM's aggregate share is highly sensitive to a calibration choice that is orthogonal to data source. This means the present test cannot, on its own, cleanly attribute Sections~\ref{sec:abm} and~\ref{sec:hazard}'s contrast to modeling paradigm rather than to data source: recalibration choice is doing comparable work to whichever of paradigm or data source the test was designed to isolate.

Both variants retain the ABM's central path-level failure: the monthly CPR path correlation stays negative at lag 0 in both variants ($-0.336$ and $-0.311$), as all estimators' are once the window's shared trend is removed (Section~\ref{sec:pathb}); Path B's distinction is level accuracy, not timing. The timing critique of Section~\ref{sec:abm} therefore survives regardless of which calibration variant is preferred, even as the dollar-recovery critique does not straightforwardly survive under the recalibrated variant.

\subsection{Cross-Design Synthesis: The Symmetric Companion Test}\label{sec:robustness-synthesis}

The symmetric companion test --- estimating the hazard framework's survival structure on a synthetic population, to check whether that structure alone recovers the benchmark independent of real data --- materially updates the cross-design picture: Path B recovers 106.0\% of the benchmark on a fully synthetic population with no real Freddie Mac data at all (96.9\% on the shared accounting basis, whose common netting applies equally here), close to its 107.0\% real-data recovery, though the raw levels correlation drops from $+0.190$ to $-0.06$ at lag 0 (a levels comparison; neither value survives detrending, Section~\ref{sec:pathb}). This decomposes the tension rather than resolving it uniformly: at the aggregate level, the hazard framework's success is not contingent on real data, which is not true of the ABM's cross-design result --- the household-choice paradigm required real covariates to reach 59.3\%, while the loan-level paradigm reaches a comparable aggregate recovery with none. At the path level, both paradigms still depend on real structure for correct timing. We treat the aggregate-level paradigm claim of Sections~\ref{sec:hazard} and~\ref{sec:conclusion} as substantially strengthened by this result, and the path-level claim as still open.

\subsection{Cross-Foundation Checks: Hybrid Pipeline and Full-Book Weighting}\label{sec:robustness-hybrid}

The hybrid pipeline referenced in Section~\ref{sec:abm-danish} holds the paper's accounting fixed while swapping the micro-foundation for loan behavior. Concretely, it runs the ABM's macro-accounting layer unchanged --- SOMA balance tracking, phased-cap netting, income-scaled curtailment, term-aware scheduled amortization, and the Danish dynamic-balance loop --- but replaces the ABM's CPR surface with Path B's microsimulated CPR path; no settlement kernel applies, since Path B settles prepayments in the month they occur. Holding the accounting fixed, the U.S.-side recovery is 97.9\% (the shared layer adds the income-scaled curtailment flow that Path B's standalone scorer omits --- \$69.56 billion over the window), and re-running the Danish leg under the Berger-recalibrated regime yields the U.S. 4.76\% / Danish 3.39\% mean-CPR pair and the $-\$99.9$ billion institutional gap that Section~\ref{sec:abm-danish} adopts as its preferred estimate. The construction matters because it removes accounting differences between frameworks as a confound of the institutional comparison: both regimes flow through one accounting layer with one empirically grounded micro-foundation.

The same layer now scores every estimator, which is what makes the recovery figures commensurable with the benchmark: the actual SOMA roll-off underlying the empirical shortfall contains borrowers' curtailment dollars, a flow no hazard simulation models, so the shared layer adds it to each simulated roll-off. The netting is identical across estimators --- \$69.56 billion for every run --- and identical \emph{by construction}, not by coincidence, with the conversion architecture stated per leg: each estimator's hazard dynamics evolve on its own internal balances (Path A drains cohort balances, Path B drains loan balances, and the Danish legs deliberately compound a dynamic counterfactual balance, Section~\ref{sec:method-abm}), but every \emph{U.S.-leg} dollar conversion is common --- simulated roll-off is rescaled month-by-month to the actual holdings path (WSHOMCB) by the book-to-sample ratio --- so the curtailment flow, computed on those same actual holdings, is a common flow rather than a per-run quantity. It is a flat 9.1-percentage-point wedge between the two bases, and we state it as such, with its scope: the invariance argument covers U.S. legs only. Danish legs compute curtailment on their own counterfactual balances inside the Danish loop; the slower Danish drain leaves higher balances and hence more curtailment dollars, a differential we quantify and print so it is checkable in the manuscript: Danish-leg curtailment is \$70.33 billion against the U.S. legs' common \$69.56 billion for the central hybrid (+\$0.77 billion), already embedded in every printed gap and immaterial to its sign or size.

On the shared basis, Path B's 107.0\% becomes 97.9\%, Path A's 119.7\% becomes 110.6\%, and the $\beta_1 = 0$ null's 97.8\% becomes 88.7\% (the +9.2-point lock-in margin is basis-invariant to within \$0.1 billion because the netting is common). The fully corrected figure comes from a joint run --- the book-coupon-reweighted simulation scored through the shared layer in one pass --- which recovers \$765.1 billion (unrounded: \$765.077 billion): 100.0\% of the \$764.7 billion benchmark at one decimal, and exactly 100.04\% as the ratio of the two unrounded figures, $765.077/764.748$, with Path A's joint run at \$894.8 billion, 117.0\%. The composed figure is conditional on the ex-ante covariate priors --- the covariate band spans 92.2\%--99.2\% on the shared basis, approximately 94.3\%--101.3\% composed (Section~\ref{sec:pathb}) --- so its precision is a property of that calibration, not a tuned result. Because the netting is common, the joint results equal the one-at-a-time deltas composed ($107.0 + 2.1 - 9.1$) exactly; we report the identity as a cross-check on the joint run, not as its source. The abstract and Sections~\ref{sec:hazard-interp} and~\ref{sec:conclusion} state recoveries on this basis.

The second check reweights the sample to the book. The Path A and Path B headline recoveries are balance-weighted estimates on the Freddie sample's composition, whose weighted-average coupon (3.9\%) sits well above the SOMA book's (2.49\%). Reweighting the simulated hazards to the book's actual coupon mix shifts Path B from 107.0\% to 109.1\% of the benchmark and Path A from 119.7\% to 126.1\% --- a composition effect of roughly two percentage points for Path B that bounds the sample-to-book extrapolation error along the coupon dimension. The agency dimension (Ginnie Mae collateral) has no internal bound and is flagged in Sections~\ref{sec:method-benchmark} and~\ref{sec:pathb}.

\subsection{Duration Quantification: Approximate Portfolio Weighted-Average Life}\label{sec:robustness-wal}

The paper's causal channel is duration extension, so we express it in maturity units rather than leaving it implicit in CPR levels. Table~\ref{tab:wal} reports the SOMA book's approximate weighted-average life (WAL) under scheduled amortization only, under the empirical CPR path, and under each estimator's implied path, at the QT window's start and end.

\begin{table}[t]
\centering
\small
\begin{threeparttable}
\caption{Approximate SOMA portfolio WAL by prepayment scenario (years).}
\label{tab:wal}
\begin{tabular}{@{}lrr@{}}
\toprule
Scenario (mean CPR) & WAL, June 2022 & WAL, Nov.\ 2025 \\
\midrule
Scheduled amortization only (0\%) & 14.7 & 12.6 \\
Empirical path (5.14\%) & 9.4 & 8.5 \\
Production ABM (11.68\%) & 6.0 & 5.6 \\
Hazard Path A (3.51\%) & 10.7 & 9.5 \\
Hazard Path B (4.76\%) & 9.7 & 8.7 \\
No-shock counterfactual, 2021 realized speeds (22.81\%) & 3.4 & 3.3 \\
\bottomrule
\end{tabular}
\begin{tablenotes}[flushleft]\footnotesize
\item Single-pool-per-cell approximation: vintage face shares from the SOMA CUSIP tabulation (2021: 43.9\%; 2020: 16.3\%; 2017--19: 6.0\%; 2022: 23.1\%; pre-2017: 10.6\%), a 90.7\%/9.1\% 30-year/15-year term split, coupon at the book's 2.49\% WAC, and each scenario's window-mean CPR held constant. Level-payment amortization with monthly SMM; November 2025 values age the same cells 42 months. Cell WAL formula: for remaining term $n$ months, $\mathrm{WAL} = \sum_{t=1}^{n} t\, p_t / (12 \sum_{t} p_t)$, with $p_t$ month-$t$ principal (scheduled level-payment principal plus $\mathrm{SMM} = 1-(1-\mathrm{CPR})^{1/12}$ of the surviving balance). Worked anchors for the scheduled-only row: a 2.49\% 30-year pool has origination WAL 16.9 years (the 2021-vintage cell, aged 12 months at June 2022, is at 16.3); the 9.1\% 15-year sleeve (origination WAL 8.0 years) and the older vintages (32.9\% of face is 2020 or earlier) pull the book blend to 14.7. Representative cell ages ($\{2022{:}\,2;\ 2021{:}\,12;\ 2020{:}\,32;\ 2017\text{--}19{:}\,60;\ \text{pre-}2017{:}\,120\}$ months at June 2022) are single-pool stand-ins recovered by matching the printed table --- a reconstruction of the original computation's assumptions, disclosed as such --- and anchored externally: ages derived from the SOMA CUSIP tabulation's own back-derived origination dates (maturity minus term, face-weighted within cells, computed with no reference to this table) run $\{2021{:}\,12;\ 2020{:}\,28;\ 2017\text{--}19{:}\,50\}$ months and give a scheduled-only blend of 14.9 years; naive vintage midpoints likewise give 14.9 (replication artifact \texttt{wal\_anchor\_results.json}). The printed 14.7 sits within 0.3 years of both alternatives. Coverage reconciliation: the 90.7\%/9.1\% split (which sums to the 99.8\% covered) is the audit parse of the \$1{,}940.9 billion book, with 0.2\% of face in other terms never modeled; the 90.6\% of Section~\ref{sec:robustness-pipeline} is the production run's own as-of-date 30-year-only coverage share --- the two parses differ by as-of date and their components round independently. Option-adjusted duration and convexity require a term-structure and prepayment-option model and are outside this paper's scope.
\end{tablenotes}
\end{threeparttable}
\end{table}

Read against the empirical column, each estimator's CPR error maps directly into effective-maturity error: the ABM's over-predicted prepayment implies a book roughly 3.4 years too short at window start, Path A's under-prediction a book 1.3 years too long, and Path B lands within 0.3 years. Relative to the no-shock world of the table's final row --- prepayments continuing at the 22.81\% mean 2021 realized CPR, implying a 3.4-year WAL --- the empirical 9.4-year WAL is a 6.0-year extension: that duration gap is what the paper's dollar benchmark prices month by month.

\subsection{Production Pipeline Fixes: Fifteen-Year MBS Fold-In and Path B Unit Correction}\label{sec:robustness-pipeline}

Two corrections made after the initial drafting of Sections~\ref{sec:abm} and~\ref{sec:hazard} change several headline figures cited throughout this paper and are documented together here for transparency, since both were caught during a routine robustness-fix pass rather than being anticipated in the original specification.

\paragraph{Fifteen-year MBS fold-in.} The empirical benchmark and ABM cohort modeling originally filtered SOMA holdings to 30-year MBS only, covering 90.6\% of face value; the remaining 15-year book (approximately 9.1\% of holdings) contributed to total portfolio size but not to cohort-correct amortization or CPR weighting. We extended the SOMA cohort parser to bucket 15-year holdings separately by coupon, never merging 15-year and 30-year buckets, and incorporated 15-year scheduled amortization into the empirical CPR back-out and the ABM's simulated roll-off. This raised coverage from 90.6\% to 99.8\% of SOMA face value and shifted the SOMA weighted-average coupon from 2.55\% to 2.49\% (Table~\ref{tab:wal}'s note reconciles these coverage components, which come from parses at different as-of dates and round independently, with the audit parse's 90.7\%/9.1\% term split).

We tested, but did not adopt, a fully native treatment in which 15-year loans also draw voluntary CPR from a 15-year-specific surface. This failed a sanity check: the ABM's payment-delta gate --- the change in \eqref{eq:annuity}'s monthly payment upon refinancing at the prevailing rate, which drives loss-aversion and mobility-desire responses elsewhere in this paper --- breaks down for seasoned 15-year borrowers, whose same-term replacement payment is nearly flat relative to their current payment. Native 15-year surfaces predicted 32\%--61\% CPR against an empirical range of roughly 5\%--8\%, enough to flip the ABM's aggregate trapped-liquidity estimate to a large negative (implausible) value. Production therefore uses a structural-only fold-in: 15-year cohorts contribute real weights and real scheduled amortization, while voluntary CPR for those cohorts continues to be drawn from the same-coupon 30-year surface. This is a scope limitation we flag rather than resolve; a 15-year-specific behavioral gate is future work.

\paragraph{Path B unit correction.} Path B's rate-gap coefficient $\beta_1$, calibrated from \citet{liebersohn2024} via the survival-function transform described in Section~\ref{sec:pathb}, is expressed as a log-hazard-ratio per 100-basis-point increase in the rate gap. An earlier implementation applied this coefficient directly to the rate gap expressed in decimal (not per-100bp) units. At a representative QT-era state --- a 3.0\% coupon against a 6.8\% market rate --- this produced a multiplier of approximately $\times$1.0026, effectively inert, versus the intended suppression of approximately $\times$0.77. (The pre-correction code carried $\beta_1$ as the raw, negative log-hazard-ratio output of \eqref{eq:beta1}'s inner transform, $\beta_1^{\mathrm{code}} = -0.0686$; multiplied by the negative decimal gap, the buggy product was positive, hence a multiplier marginally above one rather than marginally below.) Under this earlier specification, the calibrated lock-in elasticity contributed essentially nothing to Path B's simulated hazard, and the \$747 billion (97.7\%) figure it produced was generated almost entirely by the PSA baseline hazard, the involuntary-turnover floor, and the burnout term. We corrected the coefficient's application to $\exp(100\,\beta_1\,\tilde g_{i,t})$, with $\tilde g_{i,t}$ the decimal form of \eqref{eq:pathB}'s rate gap (equivalently $\exp(\beta_1 g_{i,t})$ in \eqref{eq:pathB}'s 100-basis-point units; the gap is negative when the borrower is locked in, so \eqref{eq:beta1}'s positive $\beta_1$ suppresses prepayment), consistent with Path A's convention of a positive coefficient applied to the gap in comparable units, and simultaneously corrected the Danish counterfactual's rate-gap function to rate-equivalent units so that the buyback discount properly offsets the locked-in spread. Section~\ref{sec:pathb} reports the corrected results; we surface the bug here because the \$747 billion figure appears in earlier internal drafts and in the accompanying technical documentation's version history, and a reader consulting both should understand why the two do not match.

Both corrections are structural or mechanical rather than reflecting a change in modeling philosophy, and both are additive: reproducing either fix's ``off'' state (30-year-only filtering; the original $\beta_1$ unit convention) exactly reproduces the pre-correction figures cited in earlier drafts, which we treat as a validation that the corrections are isolated to the intended code paths.

\section{Conclusion: The Securitization Trade-Off}\label{sec:conclusion}

The Federal Reserve's inability to achieve its Quantitative Tightening objectives was not solely a product of high interest rates, but a consequence of the structural design of the U.S. mortgage contract interacting with mechanisms this paper shows operate substantially outside household-level decision-making. This paper's two central empirical findings, established in Sections~\ref{sec:abm} and~\ref{sec:hazard} under the synthetic household population specified in Section~\ref{sec:method}, are complementary and stand on their own terms: the ABM --- rational, behaviorally extended, or subjected to a direct burnout test --- cannot explain the majority of the Federal Reserve's trapped liquidity, while a structurally independent hazard framework built entirely from loan-level survival dynamics, with no household utility function at all, recovers 97.9\%--110.6\% of the full empirical benchmark on the benchmark-consistent accounting basis (Section~\ref{sec:robustness-hybrid}), its microsimulation form landing within 2.1\% --- a level-accuracy distinction, not a timing one: no estimator's monthly co-movement survives detrending (Section~\ref{sec:pathb}). This contrast is, in our reading, meaningful evidence that loan-level heterogeneity, servicer behavior, and cohort-level prepayment burnout are important, and possibly dominant, mechanisms behind the Federal Reserve's stranded MBS liquidity. A separate question --- whether this contrast reflects the household-choice paradigm itself or the synthetic population it was tested on --- was taken up in Section~\ref{sec:robustness-crossdesign}, with an asymmetric answer: the hazard framework's aggregate recovery survives a fully synthetic population, the ABM's share is calibration-sensitive, and the path-level timing question remains open.

Section~\ref{sec:abm-danish}'s Danish counterfactual illustrates why this paper's revised conclusions are more modest than an earlier draft's. A mechanism-extrapolated specification initially implied a \$925.5 billion institutional gap; recalibrating against \citepos{berger2026} separately-estimated Danish elasticities collapses it by an order of magnitude, to roughly $-\$100$ billion under the hybrid pipeline, and a direct sensitivity sweep of the least-precisely-estimated input shows the remaining sign is not robustly identified (Section~\ref{sec:abm-danish}). This does not disturb Sections~\ref{sec:abm} and~\ref{sec:hazard}'s central contrast, which never depended on the Danish comparison, but it removes this paper's earlier suggestion that a Danish-style payoff rule is a ready institutional fix for the Federal Reserve's cash-flow problem specifically. The mobility relief such a rule would provide is real and well-supported by the cited literature; a corresponding institutional cash-flow benefit, or cost, is not.

A plausible reason the U.S. does not simply adopt the Danish structure is the third leg of the trade-off: the Danish market-value repurchase system is difficult to reconcile with the current U.S. fixed-income market as presently structured. The Danish system relies on strict one-to-one match funding (the Balance Principle); the United States relies on the To-Be-Announced (TBA) forward market, which requires highly homogenized, pooled mortgages to generate secondary-market liquidity. \citet{vickery2013} show that this homogeneity produces a measurable liquidity premium: TBA-eligible mortgage pools trade at significantly lower cost than otherwise comparable specified pools ineligible for TBA delivery. A market-value repurchase feature of the kind used in Denmark would introduce path-dependent, borrower-specific payoff values into each mortgage, straining the fungibility on which this premium depends. Adopting the Danish structure without a broader redesign of agency securitization rules would therefore risk sacrificing a documented, and currently uncompensated, source of secondary-market liquidity.

Ultimately, U.S. policymakers weigh three considerations, and this paper's evidence is that they do not form the strict trilemma a three-way ``trade-off'' would imply: (1) deep, homogenized secondary market liquidity (the TBA market); (2) continuous household mobility (avoiding the lock-in effect); and (3) convexity-neutral duration profiles (predictable institutional cash flows). We remain confident the U.S. has prioritized (1) at real cost to (2) --- Sections~\ref{sec:abm} and~\ref{sec:hazard} establish that household mobility is genuinely constrained and that loan-level and institutional plumbing, not household choice, governs the resulting cash-flow shortfall. What we no longer find support for is the assumption that (2) and (3) trade off sharply against each other, such that fixing mobility would cost cash flow or vice versa: Section~\ref{sec:abm-danish}'s recalibrated Danish counterfactual suggests a market-value repurchase rule could relieve (2) without materially costing (3). The more binding constraint against adopting such a rule is not a mobility--cash-flow trade-off but the TBA liquidity premium described above. The Securitization Trade-Off of this paper's title is precisely this two-sided exchange --- (1) purchased in the currency of (2) --- and this paper does not find that the price buys much of (3) either way. As long as the U.S. maintains the par-payoff MBS structure within the TBA framework, severe duration extension should be expected to recur in future tightening cycles that begin, as this one did, with a large outstanding stock of deeply in-the-money, below-market-coupon mortgages --- the state dependence documented by \citet{berger2021} and \citet{eichenbaum2022} makes the extension conditional on that stock, not on tightening per se --- and not because an alternative institutional design is unavailable, but because this paper does not identify one that clearly relieves it without cost, a conclusion more modest than, but we believe more honest than, our earlier draft's.

\subsection{Limitations and Avenues for Future Research}\label{sec:limitations}

While this paper's two estimators together substantially advance our understanding of the macroeconomic severity of the lock-in effect, and the recalibrated Danish counterfactual clarifies which part of that lock-in effect (mobility, not institutional cash flow) a market-value repurchase rule would actually relieve, several limitations remain on both sides of the framework.

Sample scale and comparability. The ABM's 10,000-household synthetic population and the hazard framework's 75,000-loan real Freddie Mac sample differ in both size and data source. Section~\ref{sec:robustness-scale} resolves the size question --- population size changes precision, not the ABM's central estimate --- and the completed cross-design pair of Section~\ref{sec:robustness-crossdesign} addresses the data-source question asymmetrically: the hazard framework's aggregate recovery is largely data-source-independent, while the ABM's is not.

Data. Freddie Mac 2017--2021 loan-level data (20 origination quarters), used for both hazard paths, is integrated via a dedicated ingest pipeline; where real data files are unavailable, results are validated first against a synthetic fixture before real-data estimation. Fifteen-year MBS, approximately 9\% of SOMA face value, remain excluded from cohort modeling on the hazard side only; on the ABM side, production folds them into cohort weighting and scheduled amortization, with voluntary CPR for those cohorts still drawn from same-coupon 30-year surfaces (Section~\ref{sec:robustness-pipeline}).

Model. On the ABM side, the monthly CPR path fit remains weak (raw $R^2$ of $-6.984$) despite a directionally sensible level response to rates, and the Danish counterfactual continues to apply U.S.-calibrated friction parameters to a population whose true institutional frictions are Danish, not American, making the institutional-gap figure an upper bound rather than a point estimate. The debt-to-income check remains front-end only, without a measure of total household debt. On the hazard side, Path A's Poisson pseudo-maximum-likelihood estimator requires a mild ridge penalty to remain well-conditioned at 295 stratum fixed effects (a stabilization device: Section~\ref{sec:patha} shows the production coefficients equal the unpenalized fit on the estimable strata), and its holdout RMSE --- approximately 38 percentage points unweighted, 2.5 points exposure-weighted --- together with a lag-0 CPR correlation of $-0.444$ more negative even than the ABM's $-0.318$, confines its evidential contribution to the aggregate dollar level, as Section~\ref{sec:patha} details; Path B's literature-calibrated hazard magnitudes, while directionally and quantitatively closer to the empirical benchmark than any ABM specification, are simulated on the 75,000-loan real Freddie Mac sample with elasticities calibrated from secondary literature rather than estimated jointly with the empirical benchmark itself, and its recovery percentage is reported across the calibration band (105.9\%--108.2\%) rather than with estimated sampling uncertainty. All headline figures are frozen in tagged run manifests in the replication repository; two freezes ran on uncommitted working trees, so manifest git hashes alone do not reproduce them --- the repository's run ledger records the exact reproduction commits.

Next steps. Several items flagged in earlier drafts of this section are now complete. The Danish-regime recalibration (Section~\ref{sec:abm-danish}) calibrates both frameworks' Danish CPR response against \citepos{berger2026} separately-estimated moving and refinancing elasticities rather than an extrapolated U.S. mechanism, collapsing the institutional gap toward zero with an unresolved sign. Path B's full 5.5\%--7.7\% elasticity band is propagated through the microsimulation (Section~\ref{sec:pathb}), which also surfaced and corrected a units error in the rate-gap coefficient's application (Section~\ref{sec:robustness-pipeline}). The symmetric companion cross-design test is complete (Section~\ref{sec:robustness-synthesis}): the hazard framework's survival structure recovers 106.0\% of the benchmark on a fully synthetic population, closely matching its real-data recovery, while its path-level timing does not transfer to synthetic data. Also complete as of this draft: Path A's calendar-month re-estimation and the concave rate-gap robustness run (Sections~\ref{sec:patha} and~\ref{sec:pathb}), the refit-and-resimulate uncertainty pass on Path A's aggregate figure (Section~\ref{sec:patha}), full-book SOMA weighting and shared-layer scoring of every estimator including the composed correction (Section~\ref{sec:robustness-hybrid}), and an external bound on the Ginnie Mae composition error (Section~\ref{sec:pathb}). Two items are bound to the pre-submission freeze as a hard gate, not left as intentions: (i) the calendar-month specification of Section~\ref{sec:patha} becomes production Path A at that freeze, with Table~\ref{tab:estimators}, the WAL table, and the refit-and-resimulate interval restated under it; (ii) every hardcoded cross-reference in the submission master converts to Word REF fields at the same freeze, and the LaTeX master must build clean in two passes with zero unresolved references or citations, the build log shipping with the submission bundle; (iii) a claims-versus-code liveness audit runs at the freeze --- every manuscript sentence presenting a mechanism as production must trace to the code path that executes it, the defect class a review round caught in the Danish counterfactual's description; and (iv) the timing-sweep generator's golden-fixture test (\texttt{tools/test\_timing\_sweep.py}) must pass at the freeze, three matcher generations having each produced a distinct silent defect. What remains open beyond that gate: adding a seasonal multiplier to Path B's baseline hazard (Section~\ref{sec:pathb}); a total-debt measure, or threshold sensitivity, for the ABM's front-end DTI gate (Section~\ref{sec:method-frictions}); and the monthly-frequency question the timing diagnosis leaves --- the empirical path's variation at that frequency is seasonal, not rate-driven, and no estimator here models it.

By testing the explanatory power of household-level utility on one side and recovering nearly the full empirical benchmark from loan-level survival dynamics on the other, Sections~\ref{sec:abm} and~\ref{sec:hazard} establish a genuine contrast under the synthetic household population this paper specifies. Section~\ref{sec:robustness}'s completed cross-design pair attributes that contrast, at the aggregate level, to modeling paradigm --- the hazard framework's recovery survives a fully synthetic population (106.0\% versus 107.0\%), while the ABM's share moves with calibration choice and data source --- and leaves the path-level timing question as the one part that remains genuinely open.

\section*{Declarations}

\textbf{Funding.} This research received no external funding.

\textbf{Conflicts of interest.} The author declares no competing interests.

\textbf{Data availability.} All inputs are drawn from publicly accessible sources: SOMA holdings data from the Federal Reserve Bank of New York; the FRED series ACTLISCOUUS, UMCSENT, DSPIC96, and related rate series from the Federal Reserve Bank of St.\ Louis; and the Freddie Mac Single-Family Loan-Level Dataset, which is available from Freddie Mac free of charge subject to its Terms of Use. Consistent with those terms, the loan-level records themselves are not redistributed with the replication package; the package instead ships the ingest pipeline that rebuilds the estimation sample from a user's own copy of the dataset, together with the cohort-level derived artifacts.

\textbf{Code availability.} The replication package --- estimation and simulation code, frozen run manifests for every headline figure (footnote~\ref{fn:manifest}), and figure-generation scripts --- accompanies this paper.

\appendix

\section{Superseded Figures and Run Ledger}\label{app:ledger}

Headline figures in this paper changed across drafts as specification errors were found and corrected. Because some superseded values appear in earlier internal drafts and in the accompanying technical documentation, we catalogue them here rather than in scattered footnotes, so that a reader consulting both always knows which figure superseded which, and why.

\begingroup\sloppy
\begin{itemize}
\item \textbf{U.S. CPR 12.3\%, empirical CPR 5.45\%, wedge 35.1pp} --- produced by a window filter that inadvertently included eight post-QT months; superseded by the bounded active-QT window (Section~\ref{sec:robustness-benchmark}).
\item \textbf{17.6pp wedge; 21--27\% Danish CPR} --- a pre-multi-cohort, rational-only calibration; superseded by the multi-vintage production specification.
\item \textbf{\$101.2B / 13.2\% ABM estimate} (\texttt{run-2026-07-04}) --- the pre-fold-in, 30-year-only tag, raw path $R^2 = -6.443$; superseded by the 15-year fold-in production tag \texttt{run-2026-07-04-15yr-foldin} (\$91.0B / 11.9\%, $R^2 = -6.984$; Section~\ref{sec:robustness-pipeline}).
\item \textbf{Danish CPR 47.1\%; wedge $\approx$35.5pp; institutional gap \$925.5B} --- the mechanism-extrapolated Danish specification; the 35.5pp wedge is the difference between its 47.1\% Danish and 11.68\% U.S. mean CPR. Superseded by the Berger-recalibrated \texttt{run-2026-07-05-berger} (Danish CPR $\approx$3.4\%; institutional gap $-\$728.4$B against that run's own \$84.5B U.S. leg; Section~\ref{sec:abm-danish}).
\item \textbf{Path B \$747B / 97.7\%} --- the pre-unit-correction microsimulation, in which the calibrated lock-in elasticity was numerically inert; superseded by \$818.5B / 107.0\% (Section~\ref{sec:robustness-pipeline}). Its corrected analogue, the deliberate $\beta_1 = 0$ null of Section~\ref{sec:hazard-interp}, recovers 97.8\%.
\item \textbf{Erratum: v15 Section~\ref{sec:robustness-hybrid} netting attribution} --- v15's Section~\ref{sec:robustness-hybrid} attributed the shared-layer netting to ``curtailment and term-aware amortization.'' False as printed: the netting is income-scaled curtailment only (\$69.56 billion); scheduled amortization is inside each estimator's own simulated roll-off. Corrected in Section~\ref{sec:robustness-hybrid}.
\item \textbf{Erratum: ``75,000 unique loans \ldots\ aggregated to 296 strata'' (v15 Table~\ref{tab:panel} note)} --- false as printed. The Path A estimation panel aggregates the \emph{full} Freddie Mac 2017--2021 origination universe (8.84 million active loans at the window's peak month); the 75,000-loan stratified sample is Path B's simulated population only. Corrected throughout this draft; population and weighting are now stated in the Table~\ref{tab:panel} note.
\item \textbf{Erratum: cross-correlation convention definition (Section~\ref{sec:pathb})} --- an earlier draft defined the lag-$k$ cross-correlation as pairing ``the empirical series at month $t$ with the simulated series at month $t+k$,'' which is the two series interchanged relative to the implementation and would have implied the opposite direction reading at the peak $k = -3$. Every direction claim in that draft (simulated trails; empirical moves first) was computed under, and verified against synthetic data on, the implemented convention and stands; the definition sentence, not the conclusions, is corrected in this draft.
\item \textbf{Erratum: \eqref{eq:beta1} printed form} --- an earlier draft printed the continuous-hazard version of the survival-function transform, $-\ln[\ln(1-(1-\delta)P_q)/\ln(1-P_q)]$, which evaluates to $0.0693$ at the central calibration; the implementation compounds discretely within the quarter ($1-P_q = (1-h_m)^3$), giving $0.0686$. Both round to the printed $\beta_1 \approx 0.069$ and both give the $\approx 0.77$ hazard ratio at QT-era gaps; the printed equation now matches the implementation. Holdout stratum-month count in Table~\ref{tab:panel} corrected from 6{,}076 to 6{,}077 (the artifact value) in the same pass.
\end{itemize}
\endgroup

Every current headline figure cites a frozen, tagged run manifest; two freezes ran on uncommitted working trees, so the repository's run ledger, not manifest git hashes, records the exact reproduction commits (Section~\ref{sec:limitations}).


\section{Path B Parameterization and Delinquency Transitions}\label{app:params}

Table~\ref{tab:params} collects every Path B parameter. All values are fixed a priori in the configuration; none is estimated on, or tuned to, the \$764.7 billion benchmark. One parameter is nonetheless measured inside the evaluation window, and benchmark-independence should not be read as window-independence: the 4\% involuntary-turnover floor is anchored to turnover observed on discount cohorts during 2023--2024 (Sections~\ref{sec:method-frictions} and~\ref{sec:pathb}), so the floor --- which sets much of the aggregate level, binding in 36\% of loan-months --- is in-sample with respect to the QT window even though it never touches the benchmark dollar figure. The covariate and default-hazard coefficients are round calibration priors. For the two prepayment covariates the signs deserve comment, because they run opposite the canonical credit-risk direction (higher FICO $\Rightarrow$ better refinancing access $\Rightarrow$ faster prepayment; higher LTV $\Rightarrow$ less equity $\Rightarrow$ slower prepayment): in the deeply out-of-the-money 2022--2025 window the refinancing channel is shut and prepayment is turnover-dominated, so credit covariates proxy mobility rather than refinancing access --- high-LTV loans are disproportionately recent purchases that churn. A Path-A-style estimation on the same stratum-month panel (Poisson PML with vintage-by-coupon fixed effects, so the credit buckets are identified from cross-bucket variation; bucket-mean covariates z-scored to the microsim pool's own moments) yields point estimates consistent with the priors' direction --- and, taken at face value, roughly double their magnitudes --- but with intervals that reject neither zero nor the opposite signs: $\hat{\beta}_F = -0.39$ (stratum-cluster bootstrap 95\% CI $[-1.51, +0.65]$) and $\hat{\beta}_L = +0.21$ ($[-0.59, +1.00]$). The estimation does not settle the sign question; it fails to contradict the priors while quantifying their sensitivity. The sensitivity brackets are correspondingly modest: a $\beta_F = \beta_L = 0$ rerun moves the recovery from 107.0\% to 108.3\%, and the panel-estimated pair moves it to 101.3\%; benchmark-independence is unaffected, since the estimation touches only the Freddie panel. The $\beta_b = 0$ rerun of Section~\ref{sec:pathb} quantifies the burnout channel the same way.

\begin{table}[t]
\centering
\footnotesize
\setlength{\tabcolsep}{4pt}
\begin{threeparttable}
\caption{Path B parameters.}
\label{tab:params}
\begin{tabular}{@{}llll@{}}
\toprule
Parameter & Value & Units & Provenance \\
\midrule
$\beta_1$ (central) & $0.069$ & per 100bp gap & \eqref{eq:beta1} at $\delta = 0.065$ \citep{liebersohn2024} \\
$\delta$ band & $0.055$--$0.077$ & quarterly decline & \citet{liebersohn2024} \\
$P_q$ & $0.06$ & quarterly prob. & turnover proxy (conversion auxiliary) \\
$\beta_F$ (FICO, z) & $-0.15$ & per SD & calibration prior; panel estimate $-0.39$ (see text) \\
$\beta_L$ (LTV, z) & $+0.10$ & per SD & calibration prior; panel estimate $+0.21$ (see text) \\
$\beta_b$ (burnout) & $-0.5$ & per unit of $B \in [0,1]$ & calibration prior \\
$\underline{h}$ (floor) & $0.0034$ & monthly SMM & 4\% annual CPR (involuntary turnover) \\
PSA speed & 100 & --- & industry seasoning convention \\
$h^{\mathrm{def}}_0$ & $3\times10^{-4}$ & monthly & calibration prior \\
$\gamma_r$ (rate stress) & $2.0$ & per decimal unit & calibration prior ($\approx$2\% per 100bp) \\
$\gamma_F$ (FICO, z) & $-0.20$ & per SD & calibration prior \\
$\gamma_L$ (LTV, z) & $+0.25$ & per SD & calibration prior \\
\bottomrule
\end{tabular}
\end{threeparttable}
\end{table}

Delinquent loans transition monthly through the matrix of Table~\ref{tab:delinq}, estimated from consecutive stratum-month \emph{modal} servicer states in the Freddie cohort panel, exposure-UPB weighted (committed as an artifact in the replication package). The division of labor with the hazard draws is one sentence: the default-hazard draw moves loans \emph{into} the pipeline (Current $\to$ D30) and terminates already-delinquent loans it hits, while the matrix governs month-to-month movement \emph{within} the pipeline --- the two never route the same loan-month. Only the delinquency rows are used in simulation, terminal states are absorbing by construction, and the forbearance state is unpopulated in the estimation panel and never entered in production runs. Because modal-state transitions are rare at the stratum-month level, the off-diagonal cells rest on thin counts, and the matrix is \emph{exposure-weighted}, so each printed probability is the exposure-UPB share of a handful of events.

We publish the full crosstab rather than leave it implicit. D30 row (21 transitions, \$58.3M row exposure): 13 cures carrying 95.6\% of the exposure weight, two D30 continuations (0.8\%), five D60 progressions (3.1\%), and one Defaulted transition (0.4\%) --- the printed 0.004 is the exposure share of that single event, whose count share would be $1/21 = 0.048$. D60 row (9 transitions, \$3.1M): six cures carrying 56.8\% and three D90+ progressions carrying 43.2\% --- the count shares are $6/9$ and $3/9$; the divergence between count and exposure shares means a few larger-balance stratum-months drive the matrix, which is exactly what an exposure-weighted estimator should do but must be said. D90+ row (38 transitions, \$76.2M): seven cures (30.3\%), one D30 (0.8\%), thirty self-continuations (68.9\%), and zero Defaulted events --- the ``---'' cells are observed zeros in a thin panel, not structural zeros, which is why serious delinquency shows no direct default exit while D30 shows one. Probabilities are printed to three decimals from these exposure shares and rows may not sum to exactly one. The matrix is accordingly reported for transparency, not precision: with cells this thin it describes the estimation panel rather than estimating a transition law, and the paper's estimates rely on the materiality bound below, not on the matrix's point values.

The thinness is consequence-free for the paper's estimates: replacing the entire matrix with either extreme (every delinquent loan cures next month; every delinquent loan defaults next month) moves trapped liquidity by \$0.39 billion --- 0.05\% of the estimate. The 482 loans delinquent at the window start are routed by the matrix from the first month (mostly cured, given the cure-heavy rows); the month-end delinquency stock thereafter never exceeds 154 of the 40,234 active loans, and delinquent UPB never exceeds 1.1\% of the pool. Both extreme-matrix runs include the initial 482, so the bound covers them.

\begin{table}[t]
\centering
\small
\begin{threeparttable}
\caption{Exposure-weighted delinquency transition matrix (modal stratum-month states).}
\label{tab:delinq}
\begin{tabular}{@{}lrrrrr@{}}
\toprule
From $\backslash$ To & Current & D30 & D60 & D90+ & Defaulted \\
\midrule
D30 & 0.956 & 0.008 & 0.031 & --- & 0.004 \\
D60 & 0.568 & --- & --- & 0.432 & --- \\
D90+ & 0.303 & 0.008 & --- & 0.689 & --- \\
\bottomrule
\end{tabular}
\begin{tablenotes}[flushleft]\footnotesize
\item Probabilities are exposure-UPB shares of thin modal-state transition counts (D30 row: 21 transitions; D60: 9; D90+: 38 --- full crosstab in the text and replication artifact). ``---'' cells are observed zeros in a thin panel, not structural zeros; rows may not sum to exactly one at printed precision. Reported for transparency, not precision: the paper's estimates rely on the \$0.39 billion extreme-matrix materiality bound in the text, not on these point values.
\end{tablenotes}
\end{threeparttable}
\end{table}

Attrition accounting: of the 75,000 sampled loans, 34,734 had already prepaid and 32 defaulted before the QT window opens, so 40,234 are active at the June 2022 start (39,752 current, 482 delinquent). The production U.S. run therefore evaluates 1,683,124 loan-months over the 42-month window --- a mean of 40,074 active loans per month, declining slightly as defaults absorb. The involuntary-turnover floor $\underline{h}$ binds in 36\% of those loan-months (Section~\ref{sec:pathb}).

\nocite{cfpb_atr,frbny_soma,fred_actliscou,fred_wshomcb,fred_dspic96,fred_umcsent}

\bibliographystyle{plainnat}
\bibliography{references}

\end{document}